% Auto-generated by analysis/build_appendix_sw.py
% Requires: \usepackage{longtable}, \usepackage{array}

\section*{Appendix: All Categorized Expert-Annotator Comments}

This appendix enumerates every expert annotator comment that participated in the S/W (strengths / weaknesses) classification of AI reviews. The full corpus is 442 substantive comments (321 item-level AI comments from \texttt{expert\_comments\_item\_level.json} plus 121 paper-level descriptive comments from \texttt{expert\_comments\_paper\_level.json}, routed via \texttt{reviewer\_unique\_items.json}). Comments labelled as being about a human reviewer, or carrying explicit \emph{item-number} references instead of free-form prose, are handled by a separate artifact and are not listed here.

Within each category, comments are sorted first by source (item-level before paper-level), then by paper id and reviewer. Each row is citable via its paper id, reviewer, and item number or paper-level slot number, so any comment can be traced back to the source JSON.

Codes have been renumbered by final frequency. $W_{16}$ (``Cannot analyze figures --- only text'') previously carried the name $W_{17}$; the old $W_{16}$ (``Aggressive / ML-style tone in non-CS fields'') is empty after manual verification and has been removed from the taxonomy.

\subsection*{Category summary}
\begin{center}
\begin{tabular}{l l r}
\hline
\textbf{Code} & \textbf{Name} & \textbf{n} \\
\hline
W1 & Missing community / field norms & 54 \\
W2 & Over-harsh / out-of-scope / unrealistic & 46 \\
W3 & Paper explicitly states X, AI says missing & 37 \\
W4 & Redundancy across the 3 AI reviewers & 28 \\
W5 & Vague / verbose / no actionable recommendation & 24 \\
W6 & Trivial / nitpicking & 16 \\
W7 & Technical term-of-art confusion & 13 \\
W8 & Cites evidence from after the preprint & 9 \\
W9 & Over-inflates small code/text inconsistencies & 9 \\
W10 & Criticizes what authors already flagged as a limitation & 6 \\
W11 & AI misreads a figure or caption & 6 \\
W12 & AI misquotes or fabricates a verbatim quote & 5 \\
W13 & AI misses supplementary content & 3 \\
W14 & Ignores authors' own cited prior work & 2 \\
W15 & AI misreads a table & 1 \\
W16 & Cannot analyze figures --- only text & 1 \\
W\_unspecified & Residual: AI judged Not Correct without specific reason & 10 \\
S1 & Statistical / methodology rigor & 45 \\
S2 & Code reading & 28 \\
S3 & Specialized niche field catch & 27 \\
S4 & Internal consistency across sections & 15 \\
S5 & Reproducibility / dependency failures & 10 \\
S6 & Big-picture / counter-narrative synthesis & 7 \\
S\_generic & Residual: AI judged Correct without specific reason & 40 \\
\hline
\textbf{Total} & & \textbf{442} \\
\hline
\end{tabular}
\end{center}

\subsection*{W1: Missing community / field norms --- n = 54}

{\footnotesize
\begin{longtable}{@{}p{3.2cm}p{11.5cm}@{}}
\toprule
\textbf{Citation} & \textbf{Expert comment} \\
\midrule
\endfirsthead
\toprule
\textbf{Citation} & \textbf{Expert comment} \\
\midrule
\endhead
\bottomrule
\endlastfoot
P2 $\cdot$ Claude 4.5 $\cdot$ item 4 $\cdot$ primary & I believe, it would make sense if this AI reviewer states what is necessary from the authors to support their claims. Yes, it is an issue that some calculations are not converging but it happens in computations/quantum chemistry. \\[3pt]
P2 $\cdot$ GPT-5.2 $\cdot$ item 2 $\cdot$ primary & While this point is essentially valid, I find it quite arguable. None AL approach provides reasonable error bounds and none ML applied to physical sciences realistically has any. AL literally helps finding new data points to improve the accuracy of models without acquiring excessive amount of them. It is typically tested empirically. Therefore, asking the authors to comment on the issue is fine but making this argument too string is unfair. \\[3pt]
P3 $\cdot$ Claude 4.5 $\cdot$ item 4 $\cdot$ primary & It is well known that 1D corrections are not working well and that is the reason why researchers using better approximations like RPMD and Instanton theory. \\[3pt]
P8 $\cdot$ GPT-5.2 $\cdot$ item 1 $\cdot$ primary & SPlot is a common technique in particle physics and widely used in many publications. \\[3pt]
P9 $\cdot$ GPT-5.2 $\cdot$ item 2 $\cdot$ primary & The comment is correct from statistical point of view, but the likelihood test cannot be done here due to how the parameters of interest are obtained. \\[3pt]
P9 $\cdot$ GPT-5.2 $\cdot$ item 5 $\cdot$ primary & It is true that these things are not documented in the paper, however they are intenally at CERN. It is not a practice to publish these as part of Nature papers so this comment is irrelevant. \\[3pt]
P9 $\cdot$ Gemini 3.0 $\cdot$ item 2 $\cdot$ primary & Full amplitude analysis is subject to a different measurement cannot be done in a CP violation measurement. \\[3pt]
P13 $\cdot$ Gemini 3.0 $\cdot$ item 4 $\cdot$ primary & Although the criticism is valid, this method has been used in many recent papers. Such maps cannot be obtained easily from many individuals, does the average maps from a sample sample are currently methodologically acceptable. \\[3pt]
P15 $\cdot$ Claude 4.5 $\cdot$ item 2 $\cdot$ primary & It's not possible to obtain intracranial data from healthy controls - it won't pass the ethical committee protocol. Unless it's a very specific population of quadriplegic patients who are getting an intracranial implant but that is extremely rare. \\[3pt]
P15 $\cdot$ Claude 4.5 $\cdot$ item 3 $\cdot$ primary & Some of this data might not be feasible to be obtained. \\[3pt]
P16 $\cdot$ GPT-5.2 $\cdot$ item 5 $\cdot$ primary & Sometimes authors release the code after publication, it could be the case for this preprint as well. \\[3pt]
P17 $\cdot$ GPT-5.2 $\cdot$ item 2 $\cdot$ primary & Although the pointing out is quite correct but I feel that it neglects that in reality the observational data is quite limited and a comprehensive data cannot be available for now. \\[3pt]
P18 $\cdot$ Claude 4.5 $\cdot$ item 2 $\cdot$ primary & The point is correct but neglecting the fact that obtaining enough observation data is always very difficult in reality. The reviewer should show some possible alternative observational data to be analyzed if proposing this claim. \\[3pt]
P18 $\cdot$ Gemini 3.0 $\cdot$ item 3 $\cdot$ primary & Neglecting the difficulty of obtaining observational data by satellites. \\[3pt]
P18 $\cdot$ Gemini 3.0 $\cdot$ item 4 $\cdot$ primary & Unfortunately, still common for physics community. \\[3pt]
P19 $\cdot$ GPT-5.2 $\cdot$ item 4 $\cdot$ primary & This comment neglects the fact that obtaining sufficient amount of observational data is impossible for astrophysical community. \\[3pt]
P20 $\cdot$ Claude 4.5 $\cdot$ item 4 $\cdot$ primary & The AI reviewer seems to lack context that it is reasonable to assume that dust attenuation can not be that high so early in the Universe from physical reasons. The AI reviewer could have provided more evidence that this can actually be a concern at these redshifts either via citations or further arguments. Also, it is reasonable to remove templates that are older than the universe, and it's unphysical to assume otherwise. \\[3pt]
P20 $\cdot$ Claude 4.5 $\cdot$ item 5 $\cdot$ primary & It's a valid point, but again the context lacks --- the preprint compares with at that time state-of-the-art measurements and the purpose is to prove the point that these brand-new observations are in striking contrast with the current best estimates with the selection function that was possible pre-JWST. \\[3pt]
P21 $\cdot$ GPT-5.2 $\cdot$ item 1 $\cdot$ primary & Valid point, but the conclusion can be tentatively excepted given the evidence in the preprint and the early/exploratory nature of the work. \\[3pt]
P22 $\cdot$ GPT-5.2 $\cdot$ item 2 $\cdot$ primary & While this point is valid, it can be acceptable to report such speculative caluculations, while clearly stating the caveats of basing this on only three objects --- which would have been accepted by a human reviewer with enough context of the field. \\[3pt]
P28 $\cdot$ Claude 4.5 $\cdot$ item 1 $\cdot$ primary & I did not do the statistical analysis the AI is claiming, but the statistical analysis in the paper is generally accepted in the community doing this type of research. \\[3pt]
P28 $\cdot$ GPT-5.2 $\cdot$ item 1 $\cdot$ primary & The main point of the critique is correct. But to the bacterial modeling community this paper does not read that it claims a universal law. However, the AI reviewer is correct with the statement ''The paper would need clearer limits on scope (e.g., ''within these assembled microcosms and this dispersal/invasion protocol'') or additional tests demonstrating the relationship persists under qualitatively different invasion settings.'' It is harsh, but correct request. \\[3pt]
P35 $\cdot$ GPT-5.2 $\cdot$ item 2 $\cdot$ primary & this is a benchmark limitation shared by all methods evaluated on BioSNAP, not a flaw unique to this paper's methodology. \\[3pt]
P37 $\cdot$ Claude 4.5 $\cdot$ item 4 $\cdot$ secondary & The authors might have not given much thought about clinical significance, but that doesn't seem like a critical problem if we treat this paper as an ML paper than a clinical paper. \\[3pt]
P37 $\cdot$ GPT-5.2 $\cdot$ item 2 $\cdot$ secondary & This AI Reviewer seems to focus on statistical rigor in research (maybe because it's following the Nature's review guideline), but I don't think the ML community these days care so much about statistical rigor any more. \\[3pt]
P41 $\cdot$ Claude 4.5 $\cdot$ item 4 $\cdot$ primary & I think the AI does understand how hard ecology is! R2 of 0.36 is impressive! \\[3pt]
P43 $\cdot$ Claude 4.5 $\cdot$ item 2 $\cdot$ primary & Privacy comes up often in these reviews. I think that it is a technical aspect that is overrepresented in the literature, but less concerning in practice than what the reviewers (AI and humans) make it seem. Especially since we are training a model on skull-stripped data. \\[3pt]
P50 $\cdot$ Claude 4.5 $\cdot$ item 2 $\cdot$ primary & Again, if this was a clinician-targeted paper, it would be important but the novelty of this paper lies in the system/engineering so normal subject is sufficient to demonstrate proof-of-concept. \\[3pt]
P50 $\cdot$ GPT-5.2 $\cdot$ item 1 $\cdot$ primary & These are important points but typically we understand that this paper is `biomedical' field paper and not `medicine' paper, meaning accuracy in clinical term is less emphasized since the novelty of this paper lies in engineering, not a clinical application. \\[3pt]
P52 $\cdot$ Claude 4.5 $\cdot$ item 2 $\cdot$ primary & This paper novelty lies in engineering/system, not a fully-deployed medical study. \\[3pt]
P53 $\cdot$ Gemini 3.0 $\cdot$ item 1 $\cdot$ primary & this comment should account for the rest of the tools available and also all the data available and how difficult is to actually obtain all taxonomy :) so the comment lack prior knowledge that an expert in the filed would know by reading other studies or having used the MSMS tools. \\[3pt]
P54 $\cdot$ Claude 4.5 $\cdot$ item 3 $\cdot$ primary & even though the criticism is on point, a human reviewer would have accounted for how difficult analysis and case studies are and that it is ok to publish as such. however yes, the authors could be careful not to overstate the application. but the potential is big and if tried in other settings than the validity will hold more. \\[3pt]
P54 $\cdot$ Gemini 3.0 $\cdot$ item 5 $\cdot$ primary & it is allowed to do such an estimation, but the programs are accounting for all of these when doing discovery and allignment. AI is probably not trained to account for this knowledge. \\[3pt]
P55 $\cdot$ Gemini 3.0 $\cdot$ item 3 $\cdot$ primary & the point 3 about the sample size, the AI is correct in what the literature suggest, but it all depends on the circumstances of the study and in this study the patient size number and the way they have been chosen is quite valuable enough. \\[3pt]
P59 $\cdot$ Claude 4.5 $\cdot$ item 1 $\cdot$ secondary & This critique applies to virtually all in vitro toxicology studies that inform policy and in my opinion this is more of a theoretical caveat. \\[3pt]
P61 $\cdot$ Gemini 3.0 $\cdot$ item 1 $\cdot$ secondary & While the concern is technically valid, it applies to virtually every study using CMIP6 for subnational health impact assessment coarse-resolution climate model output is a universal constraint. The authors partially mitigate this through the baseline calibration (Equation 5). \\[3pt]
P62 $\cdot$ Claude 4.5 $\cdot$ item 2 $\cdot$ primary & The comment is true but at this proof of concept level, there are other main problems while one can speculate that, in practice, this might not be that critical thanks to prior knowledge on samples. \\[3pt]
P62 $\cdot$ GPT-5.2 $\cdot$ item 4 $\cdot$ primary & I did not personally check the repository and the code, but the missing data could be retrieved from the authors by request. In pinciple, AI reviewer is right about the point that the paper should contain all the details and if something is missing this is a concern. However, in practice, providing everything without missing out a single point is not usually happening and this paper is already above average in terms of description of methods and details for reproducability. \\[3pt]
P62 $\cdot$ Gemini 3.0 $\cdot$ item 2 $\cdot$ primary & The comment is true but at this proof of concept level, there are other main problems while one can speculate that, in practice, this might not be that critical thanks to prior knowledge on samples. \\[3pt]
P62 $\cdot$ Gemini 3.0 $\cdot$ item 3 $\cdot$ primary & I did not personally check the repository and the code, but the missing data could be retrieved from the authors by request. \\[3pt]
P63 $\cdot$ GPT-5.2 $\cdot$ item 4 $\cdot$ primary & It is true that all the details to reproduce the results are not provided. However, the methods in the paper are well-known and there is sufficient depiction to reproduce the results, which will not be exactly the same with the paper since all the details are not provided but it is conceivable that the results would be fairly close. \\[3pt]
P65 $\cdot$ Claude 4.5 $\cdot$ item 2 $\cdot$ primary & Modified Strehl formulation is very common in metamaterial research. \\[3pt]
P65 $\cdot$ Gemini 3.0 $\cdot$ item 3 $\cdot$ primary & I don't think AI is correct. Despite having a few issues, the authors' design is one of the state-of-the-art and comparing its performance against a base design makes sense. \\[3pt]
P67 $\cdot$ Claude 4.5 $\cdot$ item 2 $\cdot$ primary & This a correct statement but here the purpose is a proof of a concept. \\[3pt]
P67 $\cdot$ Gemini 3.0 $\cdot$ item 4 $\cdot$ primary & The method has been used for nano-fabrication since early 2000. This might not be novel enough for AI, but it sure is for humans! \\[3pt]
P76 $\cdot$ Claude 4.5 $\cdot$ item 3 $\cdot$ primary & These xenograft models are typically used to assess efficacy rather than on-target, off-tumor toxicity \\[3pt]
P78 $\cdot$ Claude 4.5 $\cdot$ item 2 $\cdot$ primary & Faradaic efficiency is not highly weighed in electro-organic synthesis. AI reviewer is too harsh on this matter considering the accepted standard of the field. \\[3pt]
P78 $\cdot$ Claude 4.5 $\cdot$ item 3 $\cdot$ primary & This RPC process is assumed in this field. No need to show the machanism of the second oxidation. The AI reviewer is too harsh on this matter. \\[3pt]
P53 $\cdot$ Claude 4.5 $\cdot$ paper-level slot 2 $\cdot$ primary & Yes. the AI reviewers in general do not account to a prior knowledge of the filed so that they can make a fair point when criticizing the tool in terms of accuracy and application. Human reviewer, if chosen as a good expert in the filed, would have applied this knowledge when making the comments. This applies to all AI. \\[3pt]
P53 $\cdot$ GPT-5.2 $\cdot$ paper-level slot 1 $\cdot$ primary & Yes. the AI reviewers in general do not account to a prior knowledge of the filed so that they can make a fair point when criticizing the tool in terms of accuracy and application. Human reviewer, if chosen as a good expert in the filed, would have applied this knowledge when making the comments. This applies to all AI. AI does not pay attention to Figure commenting when it comes to introducing information in the main paper but lacks information from the supplementary files. So more connection to what info can be written in the main part and taken from supplementary is something that AI misses to account for. \\[3pt]
P53 $\cdot$ Gemini 3.0 $\cdot$ paper-level slot 3 $\cdot$ primary & Yes. the AI reviewers in general do not account to a prior knowledge of the filed so that they can make a fair point when criticizing the tool in terms of accuracy and application. Human reviewer, if chosen as a good expert in the filed, would have applied this knowledge when making the comments. This applies to all AI. \\[3pt]
P54 $\cdot$ Claude 4.5 $\cdot$ paper-level slot 3 $\cdot$ primary & AI-3 was too detailed on the statistics, the numbers of analysis which is something that the nanopore sequencing account by itself and it is not something this paper is focused on. Human reviewers that are expert in genomics would have prior knowledge about this and would usually not comment on these very specific statistics. For example the percentage of accuracy for genome alignment etc. usually you do not find a 100 \% match due to the way the softwares work. \\[3pt]
P54 $\cdot$ GPT-5.2 $\cdot$ paper-level slot 1 $\cdot$ primary & In general the AI lacks prior knowledge about what genomics community accepts as a golden standard of matching or gene alignment for some of the trained softwares. Human reviewer would know this and would not be critic about it. The AI also comments on the probable application and lacks the knowledge that even case studies of such degree can be very useful when published. \\[3pt]
P78 $\cdot$ Claude 4.5 $\cdot$ paper-level slot 1 $\cdot$ primary & AI reviewer 1 was too harsh on Faradaic economy, which is a concept that is not that highly weighed yet in synthetic organic electrochemistry. \\[3pt]
\end{longtable}
}

\subsection*{W2: Over-harsh / out-of-scope / unrealistic --- n = 46}

{\footnotesize
\begin{longtable}{@{}p{3.2cm}p{11.5cm}@{}}
\toprule
\textbf{Citation} & \textbf{Expert comment} \\
\midrule
\endfirsthead
\toprule
\textbf{Citation} & \textbf{Expert comment} \\
\midrule
\endhead
\bottomrule
\endlastfoot
P1 $\cdot$ Claude 4.5 $\cdot$ item 3 $\cdot$ primary & Overall, I agree with this comment. However, I would not claim that what the authors do is completely wrong. One may simply ask for doing that additional set of experiments. \\[3pt]
P2 $\cdot$ Claude 4.5 $\cdot$ item 2 $\cdot$ primary & However, requesting CBS may be too much. \\[3pt]
P4 $\cdot$ Gemini 3.0 $\cdot$ item 1 $\cdot$ primary & It is an architectural choice and it is valid to have advantages by using it. \\[3pt]
P6 $\cdot$ Claude 4.5 $\cdot$ item 1 $\cdot$ primary & This is the least interesting aspect of the paper. None of the 3 human reviewers even bothered with it, because the paper is using safety as an excuse to not talk about the MR engineering. The AI may be right (I am not a statistics expert), but it is a marginal issue. \\[3pt]
P6 $\cdot$ Claude 4.5 $\cdot$ item 2 $\cdot$ primary & The entire safety portion of the paper is just filler to justify not talking about the MR engineering. This is partially the paper's fault, but the human reviewers were not fooled by the misdirection. \\[3pt]
P6 $\cdot$ Claude 4.5 $\cdot$ item 4 $\cdot$ primary & Again irrelevant. The safety aspects belong in another paper and the duration of the follow-up is very field-specific. While it is true that the Grant et al. study looked at longer time horizons, I doubt that a follow-up of 1 week vs. 2-4 weeks makes any difference. \\[3pt]
P6 $\cdot$ GPT-5.2 $\cdot$ item 1 $\cdot$ secondary & not focus of the paper \\[3pt]
P8 $\cdot$ Claude 4.5 $\cdot$ item 2 $\cdot$ primary & This is not relevant and cannot be covered in this paper. \\[3pt]
P8 $\cdot$ GPT-5.2 $\cdot$ item 4 $\cdot$ primary & It's true that the assumption is made, however these inputs are not subject to a measurement in this paper. \\[3pt]
P28 $\cdot$ Claude 4.5 $\cdot$ item 3 $\cdot$ primary & The statement of AI is correct, however fitting models against experimental data is beyond the scope of the paper. \\[3pt]
P28 $\cdot$ Claude 4.5 $\cdot$ item 4 $\cdot$ primary & The statements are correct, but too harsh and are beyond the scope of the paper. \\[3pt]
P28 $\cdot$ Claude 4.5 $\cdot$ item 5 $\cdot$ primary & The statement is correct, but a threshold must be chosen so the comment, although technically correct, is to harsh for a review. \\[3pt]
P28 $\cdot$ GPT-5.2 $\cdot$ item 2 $\cdot$ primary & This is too harsh critique. The authors do not claim that their result is universal, although the AI would have had a good point if they did. This critique is somewhat along the lines of Human Reviewer 2. \\[3pt]
P28 $\cdot$ GPT-5.2 $\cdot$ item 3 $\cdot$ primary & The AI is finding problems with external packages and software. This is too much, I did not go down the rabbit hole to evaluate it. \\[3pt]
P28 $\cdot$ Gemini 3.0 $\cdot$ item 3 $\cdot$ primary & The AI here is in principle correct. However, this critique is too harsh. From the context it is clear that the scope of the paper is somewhat limited. This critique is similar to Human Reviewer 2. \\[3pt]
P29 $\cdot$ Claude 4.5 $\cdot$ item 1 $\cdot$ primary & This is technically correct, but outside the scope of the paper. \\[3pt]
P29 $\cdot$ Claude 4.5 $\cdot$ item 4 $\cdot$ primary & Correct, but out of the scope of the paper. \\[3pt]
P30 $\cdot$ Claude 4.5 $\cdot$ item 2 $\cdot$ primary & Out of the scope of the paper. \\[3pt]
P30 $\cdot$ Gemini 3.0 $\cdot$ item 5 $\cdot$ primary & Out of the scope of the paper. \\[3pt]
P31 $\cdot$ Claude 4.5 $\cdot$ item 4 $\cdot$ primary & Too harsh critique. \\[3pt]
P37 $\cdot$ Claude 4.5 $\cdot$ item 2 $\cdot$ secondary & The first two quote-comment pairs are logically valid, but I don't think they are enough to reject the paper. However, the third pair is more critical, as one of the core claims (i.e. the importance of date embedding) is negated by the empirical result. \\[3pt]
P42 $\cdot$ Claude 4.5 $\cdot$ item 3 $\cdot$ primary & Technically the AI is correct, but I think much like the other paper I reviewed, the AI does not understand the complexity of studying ecosystems and so has unrealistic expectations of the resultant classification performance. The signal in the models is strong even with the moderate classification performance. \\[3pt]
P42 $\cdot$ GPT-5.2 $\cdot$ item 3 $\cdot$ primary & While technically correct again, I think the AI does not understand the complexity and challenges of studying ecosystems; however, there is some causal language in the paper that could be revised. \\[3pt]
P43 $\cdot$ Claude 4.5 $\cdot$ item 1 $\cdot$ primary & The interesting part about the paper is more about the model itself rather than the federation. \\[3pt]
P46 $\cdot$ Claude 4.5 $\cdot$ item 2 $\cdot$ primary & Yes, but what can you do?! \\[3pt]
P46 $\cdot$ Claude 4.5 $\cdot$ item 4 $\cdot$ primary & Yes, but it is a valid and good comparison! \\[3pt]
P46 $\cdot$ GPT-5.2 $\cdot$ item 3 $\cdot$ primary & It is correct that this is the case, but PG cannot be effectively implemented on the GPU, whereas polyBERT can. A comment in the text may be good, but to me this is not a big deal. \\[3pt]
P49 $\cdot$ Claude 4.5 $\cdot$ item 2 $\cdot$ secondary & This critic is factually accurate but it highlights a limitation that the authors cannot realistically defend against or resolve by modifying their data, making it a non-constructive critique. \\[3pt]
P52 $\cdot$ Claude 4.5 $\cdot$ item 4 $\cdot$ primary & Given this is academic paper, it's difficult to ask authors to provide full biocompatibility data which would take too long. \\[3pt]
P54 $\cdot$ Claude 4.5 $\cdot$ item 1 $\cdot$ primary & Bit too rigid criticizm however, it does point out when authors overstate their point trying to sell their paper. \\[3pt]
P59 $\cdot$ Claude 4.5 $\cdot$ item 4 $\cdot$ secondary & The dose response non-linearity for total PM 2.5 mass has been characterized, but there is no established dose-response function for toxic potency-adjusted PM2.5. This is a future research direction, not something the authors could realistically incorporate. Also citation 5, the link is wrong. \\[3pt]
P59 $\cdot$ GPT-5.2 $\cdot$ item 4 $\cdot$ secondary & Applies an unrealistically stringent reproducibility standard. This standard of reproducibility, while ideal, is rarely met in chemical transport modeling studies published in top journals. \\[3pt]
P59 $\cdot$ Gemini 3.0 $\cdot$ item 3 $\cdot$ secondary & reviewer seems to recognize that this is not a methodological flaw - ''it does not necessarily invalidate the relative ranking'' but still frames it as a significant concern. It can be a valuable addition to the discussions but its not affecting the conclusions. \\[3pt]
P62 $\cdot$ Claude 4.5 $\cdot$ item 3 $\cdot$ primary & The comment is true but one can easily make an counter-argument of these not being a huge problem in controlled experiments. \\[3pt]
P65 $\cdot$ Claude 4.5 $\cdot$ item 1 $\cdot$ primary & AI got this wrong. Comparing simulation and experimental results is not fair. \\[3pt]
P65 $\cdot$ Gemini 3.0 $\cdot$ item 5 $\cdot$ primary & It looks like this was a badly trained AI platform. What could the authors do if their imaging takes place at certain wavelength values other than averaging the discrete outcome? This was my first study on Human vs AI reviewers. It is interesting to see that one of three AI platforms can come up with a totally different review than the other two. Clearly its training was done with different sources than the other two. LTSM was an issue 5-10 years ago. The AI reviewer III shows me that it might be still an issue for some AI platforms. There is a more significant alignment between AI reviewers I and II. THey saw some important points which were not caught by human reviewers. \\[3pt]
P66 $\cdot$ GPT-5.2 $\cdot$ item 1 $\cdot$ primary & This is such a harsh comment for an experimental study. The authors provided all the datasets. The measurements were taken several times and averages were reported. The real world measurements are not perfect or noise free. What else could authors do? \\[3pt]
P66 $\cdot$ GPT-5.2 $\cdot$ item 2 $\cdot$ primary & I again do not agree with the AI reviewer. The wavelength is close to 800 nm and the tip is 20-25 nm wide, which is 1/30th - 1/40th of the wavelength. We can treat it as a point dipole. The AI reviewer's comment is too harsh! \\[3pt]
P67 $\cdot$ GPT-5.2 $\cdot$ item 4 $\cdot$ primary & I understand AI's request, which is fair. But this shouldn't undermine the authors' contribution \\[3pt]
P68 $\cdot$ Gemini 3.0 $\cdot$ item 1 $\cdot$ primary & This is technically correct, but morphological traits are very correlated, which makes the underlying dimensionality much lower than the number of coordinates. So, I am two ways with regards to this comment. On one way, the criticism is correct. On the other, it is likely overstated. \\[3pt]
P72 $\cdot$ Claude 4.5 $\cdot$ item 5 $\cdot$ primary & DeBreak is more computationally efficient than one of the three standard callers it is compared to (pbsv), so I think this criticism is overblown. \\[3pt]
P72 $\cdot$ Gemini 3.0 $\cdot$ item 4 $\cdot$ primary & One could always compare to more tools, but I probably would not have asked the authors to include an additional tool unless there was a specific reason for including it. e.g. one of the human reviewers' requested a comparison to PBHoney and had a good reason for doing so. \\[3pt]
P77 $\cdot$ Claude 4.5 $\cdot$ item 5 $\cdot$ secondary & A definitive answer to this question may not be readily available. Nevertheless, it would be valuable if the authors could provide insight into the different reactivities of various metals. \\[3pt]
P6 $\cdot$ Claude 4.5 $\cdot$ paper-level slot 1 $\cdot$ primary & AI Reviewer 1 focused too much on the safety aspects of the paper. 70\% of the review was about safety, whereas the paper is primarily about MRI methodology. I am sure that some of the statistics observations were original, but they were irrelevant in my opinion. \\[3pt]
P8 $\cdot$ Claude 4.5 $\cdot$ paper-level slot 1 $\cdot$ primary & Human review 2 is not available in the checklist. I noticed that there was one question in the AI reviewer that was suggesting to provide a measurements of some parameters which is out of the scope of the analysis done there and also some trivial questions on splot that are not applicable. That would never be asked by a human reviewer. In the AI reviewer 2 there were some trivial questions that should not be asked in Nature paper \\[3pt]
P28 $\cdot$ GPT-5.2 $\cdot$ paper-level slot 1 $\cdot$ primary & Yes, the AI reviewer 1 was talking about the equations and math, but those comments are relevant for external software, not authored in this paper. \\[3pt]
\end{longtable}
}

\subsection*{W3: Paper explicitly states X, AI says missing --- n = 37}

{\footnotesize
\begin{longtable}{@{}p{3.2cm}p{11.5cm}@{}}
\toprule
\textbf{Citation} & \textbf{Expert comment} \\
\midrule
\endfirsthead
\toprule
\textbf{Citation} & \textbf{Expert comment} \\
\midrule
\endhead
\bottomrule
\endlastfoot
P1 $\cdot$ Gemini 3.0 $\cdot$ item 1 $\cdot$ primary & The authors write explicitly that they employ finite-size correction. \\[3pt]
P3 $\cdot$ Claude 4.5 $\cdot$ item 1 $\cdot$ primary & I guess one can criticize it but authors have already provided a reasonable evaluation. \\[3pt]
P3 $\cdot$ Claude 4.5 $\cdot$ item 3 $\cdot$ primary & They give values they got in their experiments which is appropriate for this case. \\[3pt]
P4 $\cdot$ GPT-5.2 $\cdot$ item 3 $\cdot$ primary & The authors provide the code and it does not matter whether it is on GitHub or GitLab. \\[3pt]
P4 $\cdot$ Gemini 3.0 $\cdot$ item 5 $\cdot$ primary & The code is available. \\[3pt]
P5 $\cdot$ Claude 4.5 $\cdot$ item 4 $\cdot$ primary & Formally, this criticism is correct, but it can be omitted because the authors provide molecular dynamics simulations for other benchmark systems. \\[3pt]
P6 $\cdot$ Gemini 3.0 $\cdot$ item 1 $\cdot$ primary & The AI made factual errors. The sequence parameters were reported, and the methods section starts with this. It is true that the MRI methodology does not have a subheading, which is why the AI got confused. \\[3pt]
P8 $\cdot$ Claude 4.5 $\cdot$ item 1 $\cdot$ primary & The comment make sense however the authors consider more complicated model as a systematic study. \\[3pt]
P9 $\cdot$ GPT-5.2 $\cdot$ item 4 $\cdot$ primary & The comment is too vague of what is insufficient in the paper since there is some discussion on the resonant part. \\[3pt]
P12 $\cdot$ Gemini 3.0 $\cdot$ item 3 $\cdot$ secondary & The concern about J1 being a non-detection is valid, but the Lorentz factor is derived including J2. \\[3pt]
P17 $\cdot$ Gemini 3.0 $\cdot$ item 1 $\cdot$ primary & The pointing out on how to define classification during non-equilibrium process is appropriate but it neglects that the paper considers a response to linear perturbation and a near equilibrium initial profiles of observation and PIC simulations, leading to inadequately strong criticism to the authors' claim. The wording seems common on ML community unfriendly review comment but less common in physics community review. \\[3pt]
P18 $\cdot$ GPT-5.2 $\cdot$ item 2 $\cdot$ primary & AI confused by a confusing authors' explanation, mixing 1D analyses with 2D reconnection case without explicitly clarifying. Just suggesting clarifying the explanation on that part is more appropriate as a review comment. \\[3pt]
P25 $\cdot$ Gemini 3.0 $\cdot$ item 3 $\cdot$ primary & I was doubting between correct and not correct. What the criticism reffers to are instantaneous correlations, due to the wave there are lagged correlations. \\[3pt]
P28 $\cdot$ Gemini 3.0 $\cdot$ item 1 $\cdot$ primary & The paper claims usage of a Matlab method, and does not explicitly state the equations mentioned in the review here. This is clearly a fail of the AI reviewer. \\[3pt]
P30 $\cdot$ Claude 4.5 $\cdot$ item 1 $\cdot$ primary & It is clear from the context that only QS systems are included. \\[3pt]
P30 $\cdot$ GPT-5.2 $\cdot$ item 1 $\cdot$ primary & It is clear from the context that QS systems only are considered. \\[3pt]
P31 $\cdot$ Gemini 3.0 $\cdot$ item 1 $\cdot$ primary & Correct, but explained in the paper that this is inevitable. \\[3pt]
P34 $\cdot$ Gemini 3.0 $\cdot$ item 1 $\cdot$ primary & The test should be done \textasciitilde{}3000 times consistently with the paper, single run shouldn't be interpreted like this. \\[3pt]
P37 $\cdot$ GPT-5.2 $\cdot$ item 4 $\cdot$ primary & the model architecture for finetuning is provided (ICDBartForSequenceBinaryClassification in icdmodelbart.py) but the training script is not. \\[3pt]
P43 $\cdot$ Claude 4.5 $\cdot$ item 3 $\cdot$ primary & This comment disregards the existence of out-of-sample sites \\[3pt]
P44 $\cdot$ Claude 4.5 $\cdot$ item 2 $\cdot$ primary & It's failing to take into account the reduced set, which naturally leads to higher p-values. \\[3pt]
P50 $\cdot$ Gemini 3.0 $\cdot$ item 3 $\cdot$ primary & There is actually `methods' text in the manuscript. It probably wasn't properly inputed into AI. \\[3pt]
P55 $\cdot$ Claude 4.5 $\cdot$ item 2 $\cdot$ primary & the age is defined as above 18 so there is a need for more info to make such a comment \\[3pt]
P59 $\cdot$ Gemini 3.0 $\cdot$ item 1 $\cdot$ secondary & This critic assumed ''zero-out'' meant removing all emissions (PM, SO2) from each sector but in the paper authors clearly mentioned that they zero out only primary PM 2.5 and they akcnowldege this in the limitations section as well (i.e ignoring secondary aerosol contribution). \\[3pt]
P61 $\cdot$ Claude 4.5 $\cdot$ item 5 $\cdot$ secondary & this is factually incorrect, the paper explicitly addresses this issue. Lines 489--496 describe a calibration procedure (Equation 5) that adjusts modeled PM2.5 concentrations against observed baseline period data (2012--2017). \\[3pt]
P61 $\cdot$ Gemini 3.0 $\cdot$ item 4 $\cdot$ secondary & The paper's title, abstract, and methods explicitly scope the study to PM2.5. The reviewer's argument that the paper's broader framing around ''air pollution-related deaths'' and SDGs implicitly requires ozone consideration is a stretch no single study is obligated to cover all pollutants, and the PM2.5 focus is clearly stated throughout. \\[3pt]
P61 $\cdot$ Gemini 3.0 $\cdot$ item 5 $\cdot$ secondary & The paper actually uses the GBD2019 MR-BRT model (line 470), not the older IER function, and explicitly discusses the GEMM comparison (lines 547--550), acknowledging that GEMM yields higher estimates. \\[3pt]
P62 $\cdot$ GPT-5.2 $\cdot$ item 2 $\cdot$ secondary & Marked it maginal, because important metrics are provided in their experimental context. \\[3pt]
P66 $\cdot$ Claude 4.5 $\cdot$ item 2 $\cdot$ primary & The authors mention repetition of measurements and using averaged results in the end. I don't think this is a fair comment. \\[3pt]
P66 $\cdot$ Claude 4.5 $\cdot$ item 3 $\cdot$ primary & The authors mention that their resolution is limited by the width of their tip and their experimental results in a good agreement with theoretical results (10 nm). There is no issue here! \\[3pt]
P67 $\cdot$ GPT-5.2 $\cdot$ item 3 $\cdot$ primary & The zip file I downloaded has all the *.mat files \\[3pt]
P67 $\cdot$ GPT-5.2 $\cdot$ item 3 $\cdot$ secondary & In the code bundle that was provided to me, the .mat files were empty. I believe this is what the AI reviewer is referring to. However, I could find the actual .mat files on GitHub. \\[3pt]
P67 $\cdot$ Gemini 3.0 $\cdot$ item 3 $\cdot$ primary & They wrote ''to simulate the process.'' \\[3pt]
P67 $\cdot$ Gemini 3.0 $\cdot$ item 5 $\cdot$ primary & They already showed field enhancement far from the contacts. \\[3pt]
P74 $\cdot$ Claude 4.5 $\cdot$ item 3 $\cdot$ primary & The point is valid but overblown given that the authors also assess accuracy on non-human genomes. \\[3pt]
P74 $\cdot$ GPT-5.2 $\cdot$ item 2 $\cdot$ primary & Another valid point that is somewhat offset by the benchmarking on other organisms. \\[3pt]
P77 $\cdot$ Claude 4.5 $\cdot$ item 2 $\cdot$ primary & The substrates described in this paper do not have this issue. \\[3pt]
\end{longtable}
}

\subsection*{W4: Redundancy across the 3 AI reviewers --- n = 28}

{\footnotesize
\begin{longtable}{@{}p{3.2cm}p{11.5cm}@{}}
\toprule
\textbf{Citation} & \textbf{Expert comment} \\
\midrule
\endfirsthead
\toprule
\textbf{Citation} & \textbf{Expert comment} \\
\midrule
\endhead
\bottomrule
\endlastfoot
P5 $\cdot$ Claude 4.5 $\cdot$ item 2 $\cdot$ primary & To a large extent, all three AI reviewers raise very similar points. Thus, the overall review would not cover many substantial points human reviewer raised. \\[3pt]
P42 $\cdot$ GPT-5.2 $\cdot$ item 1 $\cdot$ primary & Same comment as AI 2 item 2. \\[3pt]
P46 $\cdot$ Gemini 3.0 $\cdot$ item 1 $\cdot$ primary & Same as point 1 of other reviewer. \\[3pt]
P59 $\cdot$ GPT-5.2 $\cdot$ item 2 $\cdot$ secondary & This item largely duplicates AI Review \#1 Item 1. Despite being more elaborately argued, it doesn't add a fundamentally new dimension. The authors should temper their language from ''health-oriented'' to something like ''toxicity-informed'' and explicitly acknowledge that TPAE is an exploratory index rather than a calibrated health impact function is easy to implement. \\[3pt]
P59 $\cdot$ GPT-5.2 $\cdot$ item 3 $\cdot$ secondary & Essentially the same critique as AI Review \#2 Item 1. It actually creates a conceptual tension: secondary sulfate/nitrate are widely considered less toxic, so including them would increase each sector's attributable concentration while simultaneously lowering the average toxicity per unit mass. The reviewer asserts this ''can materially change'' rankings but provides no quantitative basis for that claim. \\[3pt]
P60 $\cdot$ Claude 4.5 $\cdot$ item 1 $\cdot$ secondary & This is a well-constructed critique but largely redundant with the AI reviewer1's Items 2 and 3. The main added value is the specific Ali et al. (2025) reference demonstrating synergistic effects with quantified RRs, which provides stronger empirical backing than the first reviewer offered for the same point. \\[3pt]
P60 $\cdot$ Claude 4.5 $\cdot$ item 2 $\cdot$ secondary & This is essentially the same critique as the AI reviewer 1 Item 2, which already identified that removing COVID-19 deaths doesn't make the remaining mortality independent of pandemic context, and that the model trained on primarily pre-pandemic data fails to capture altered vulnerability. \\[3pt]
P60 $\cdot$ GPT-5.2 $\cdot$ item 1 $\cdot$ secondary & AI Reviewer 1 (Items 2 and 3) identified both the violated independence assumption and the misleading ''compound'' terminology. AI Reviewer 2 (Item 1) restated it with additional references. This reviewer articulates it more formally mentioning joint models, attributable fraction partitioning, and counterfactual decomposition but the core insight is identical \\[3pt]
P60 $\cdot$ GPT-5.2 $\cdot$ item 2 $\cdot$ secondary & This is a thorough and well-referenced elaboration of a point other reviwers already made more concisely. The added value is the explicit connection to the Vicedo-Cabrera projection framework, which helps explain why the authors may have made this choice (importing a projection-style baseline into a retrospective analysis) and why it's inappropriate here. \\[3pt]
P61 $\cdot$ Gemini 3.0 $\cdot$ item 2 $\cdot$ secondary & This item is essentially a less developed version of what AI Review \#2 Item 2 covers more comprehensively. \\[3pt]
P62 $\cdot$ Claude 4.5 $\cdot$ item 3 $\cdot$ secondary & Aligned with Item 7 of Human Reviewer 3. \\[3pt]
P62 $\cdot$ GPT-5.2 $\cdot$ item 1 $\cdot$ secondary & Aligned with Item 7 of Human Reviewer 3, Item 3 of AI Reviewer 1, and Item 1 of AI Reviewer 2. But not addressing other important challenges in in vivo settings like Human Reviewer 3. \\[3pt]
P62 $\cdot$ GPT-5.2 $\cdot$ item 3 $\cdot$ primary & Again the same comment about NMF rank that is common in all AI reviews. \\[3pt]
P62 $\cdot$ GPT-5.2 $\cdot$ item 3 $\cdot$ secondary & Aligned with Item 2 of AI Reviewer 1 and Item 2 of AI Reviewer 2. \\[3pt]
P62 $\cdot$ GPT-5.2 $\cdot$ item 4 $\cdot$ secondary & Aligned with Item 3 of AI reviewer 2. Marked marginal by the same reason. \\[3pt]
P62 $\cdot$ Gemini 3.0 $\cdot$ item 1 $\cdot$ secondary & Aligned with Item 7 of Human Reviewer 3, and Item 3 of AI Reviewer 1. But not addressing other important challenges in in vivo settings like Human Reviewer 3. \\[3pt]
P62 $\cdot$ Gemini 3.0 $\cdot$ item 2 $\cdot$ secondary & Aligned with Item 2 of AI Reviewer 1. \\[3pt]
P72 $\cdot$ GPT-5.2 $\cdot$ item 3 $\cdot$ primary & A fair point, echoing a comment from one of the human reviewers. \\[3pt]
P72 $\cdot$ GPT-5.2 $\cdot$ item 4 $\cdot$ primary & Echoes a point by one of the human reviewers that the authors could emphasize that their benchmarking results do not encompass all SV types handled by DeBreak. \\[3pt]
P4 $\cdot$ Claude 4.5 $\cdot$ paper-level slot 3 $\cdot$ primary & I found the point on the comparison with prior work important. Other points are rather similar to other AI reviewers. \\[3pt]
P22 $\cdot$ Claude 4.5 $\cdot$ paper-level slot 3 $\cdot$ primary & In general, they pick up on the same major points \\[3pt]
P22 $\cdot$ GPT-5.2 $\cdot$ paper-level slot 1 $\cdot$ primary & In general, they pick up on the same major points \\[3pt]
P22 $\cdot$ Gemini 3.0 $\cdot$ paper-level slot 2 $\cdot$ primary & In general, they pick up on the same major points \\[3pt]
P62 $\cdot$ Claude 4.5 $\cdot$ paper-level slot 1 $\cdot$ primary & None. In this paper's specific case, all human reviewers were exceptional and their combined feedback covered every possible aspect. All AI reviewers to some extent failed to give feedback that could possibly improve the manuscript while keeping the intricacies and difficulties of optics+bio experimental work landscape. \\[3pt]
P62 $\cdot$ GPT-5.2 $\cdot$ paper-level slot 3 $\cdot$ primary & None. In this paper's specific case, all human reviewers were exceptional and their combined feedback covered every possible aspect. All AI reviewers to some extent failed to give feedback that could possibly improve the manuscript while keeping the intricacies and difficulties of optics+bio experimental work landscape. \\[3pt]
P62 $\cdot$ Gemini 3.0 $\cdot$ paper-level slot 2 $\cdot$ primary & None. In this paper's specific case, all human reviewers were exceptional and their combined feedback covered every possible aspect. All AI reviewers to some extent failed to give feedback that could possibly improve the manuscript while keeping the intricacies and difficulties of optics+bio experimental work landscape. \\[3pt]
P67 $\cdot$ Claude 4.5 $\cdot$ paper-level slot 3 $\cdot$ secondary & The review of this AI is similar to the review of AI reviewer 1 but I find that the AI reviewer 1 commented on matters that are more important first. \\[3pt]
P72 $\cdot$ GPT-5.2 $\cdot$ paper-level slot 3 $\cdot$ primary & No. This AI Reviewer was the only one to complain about the assembly-based callset being treated as ground-truth, but at least one other reviewer was clearly aware that this was not the authors' true intention---they were just using it as a decently accurate set that might be worth comparing alignment-based methods to. \\[3pt]
\end{longtable}
}

\subsection*{W5: Vague / verbose / no actionable recommendation --- n = 24}

{\footnotesize
\begin{longtable}{@{}p{3.2cm}p{11.5cm}@{}}
\toprule
\textbf{Citation} & \textbf{Expert comment} \\
\midrule
\endfirsthead
\toprule
\textbf{Citation} & \textbf{Expert comment} \\
\midrule
\endhead
\bottomrule
\endlastfoot
P1 $\cdot$ Claude 4.5 $\cdot$ item 2 $\cdot$ primary & The valid point is that authors refer to ref. 21 as providing the experimental value. However, ref. 21 takes the corresponding reference value from a different work. This could be requested to correct. But, in general, the comment/critique by this reviewer is rather pointless. \\[3pt]
P5 $\cdot$ GPT-5.2 $\cdot$ item 1 $\cdot$ primary & While I agree with the points raised by this AI reviewer, I found them somewhat too lengthy and repetitive. \\[3pt]
P5 $\cdot$ GPT-5.2 $\cdot$ item 4 $\cdot$ primary & Similar to the previous item, this AI reviewer tends to overemphasize each quotation and comment. It would be sufficient to state that repeating the simulation with different starting points and random seeds is necessary to estimate the variance and thus draw more reliable conclusions. This AI reviewer also appears to misuse the term of ''uncertainty''. \\[3pt]
P6 $\cdot$ GPT-5.2 $\cdot$ item 1 $\cdot$ primary & Too verbose. I am not a statistics expert, but this is a paper about MR methodology. The first comment should not be six paragraphs about biological effects. \\[3pt]
P6 $\cdot$ GPT-5.2 $\cdot$ item 2 $\cdot$ primary & Too verbose. The previous two AI reviewers did a much better job pinpointing the main issues. \\[3pt]
P17 $\cdot$ Gemini 3.0 $\cdot$ item 2 $\cdot$ primary & The pointing out is basically correct but needs more comments. Maybe a more appropriate review comment on this point is to suggest to take more effort on validating why the authors used MMS observation data with reconnection to compare 1D non-reconnection current sheet PIC simulations. Basically, it criticized too much. \\[3pt]
P22 $\cdot$ GPT-5.2 $\cdot$ item 1 $\cdot$ primary & Thorough and well articulated, but perhaps unecessarily long. \\[3pt]
P22 $\cdot$ GPT-5.2 $\cdot$ item 3 $\cdot$ primary & Too long. \\[3pt]
P29 $\cdot$ Claude 4.5 $\cdot$ item 2 $\cdot$ primary & This critique, technically correct, is confusing the significance of some of the results. \\[3pt]
P29 $\cdot$ GPT-5.2 $\cdot$ item 2 $\cdot$ primary & There are some correct elements here, but it is stated in a super confusing way. \\[3pt]
P29 $\cdot$ GPT-5.2 $\cdot$ item 3 $\cdot$ primary & This has an element of truth, but evidence is not provided well, and the statement is convoluted. \\[3pt]
P29 $\cdot$ Gemini 3.0 $\cdot$ item 3 $\cdot$ primary & This might be correct but it is stated in a convoluted way. \\[3pt]
P36 $\cdot$ Gemini 3.0 $\cdot$ item 3 $\cdot$ primary & The evidence consists only of paper quotes restated with the reviewer's interpretation, without external citations or concrete reasoning to support the reverse causality claim \\[3pt]
P41 $\cdot$ Claude 4.5 $\cdot$ item 1 $\cdot$ primary & As with other reviews, it is unclear what the recommendation for addressing this criticism is. \\[3pt]
P41 $\cdot$ Claude 4.5 $\cdot$ item 3 $\cdot$ primary & Similar to AI review \#1, this comment is true and yet the signal in the results is still strong, it's unclear from the comment whether this is a fundamental criticism or something that requires addressing in the discussion. \\[3pt]
P41 $\cdot$ GPT-5.2 $\cdot$ item 1 $\cdot$ primary & As with other reviews, it is unclear what the recommendation for addressing this criticism is. \\[3pt]
P41 $\cdot$ Gemini 3.0 $\cdot$ item 1 $\cdot$ primary & The review does not provide any suggestions for how to address this criticism. \\[3pt]
P41 $\cdot$ Gemini 3.0 $\cdot$ item 2 $\cdot$ primary & It is unclear from the criticism how this issue results in bias per se, which is a much greater concern than ``noise'', given the strong signal in the results of the paper. \\[3pt]
P49 $\cdot$ GPT-5.2 $\cdot$ item 5 $\cdot$ secondary & While this AI reviewer (and others) is excellent at pointing out flaws, it fails to provide constructive alternatives or actionable suggestions on how the authors should address these criticisms to improve the paper. \\[3pt]
P60 $\cdot$ GPT-5.2 $\cdot$ item 4 $\cdot$ secondary & The point is valid but the level of detail suggests the reviewer may be optimizing for thoroughness over insight. \\[3pt]
P77 $\cdot$ GPT-5.2 $\cdot$ item 4 $\cdot$ secondary & I could not understand the reviewer's main point. \\[3pt]
P78 $\cdot$ Gemini 3.0 $\cdot$ item 2 $\cdot$ secondary & Although the point mentioned by the reviewer is valid, more evidence (references) should be given. \\[3pt]
P5 $\cdot$ GPT-5.2 $\cdot$ paper-level slot 2 $\cdot$ primary & I found the items themselves to be reasonable, but the underlying argumentation was diffuse and overcrowded with statements that were largely irrelevant. Overall, as noted in the checklist, all reviewers raised very similar points. However, the first reviewer provided the strongest arguments, whereas this reviewer's arguments were the weakest. \\[3pt]
P6 $\cdot$ GPT-5.2 $\cdot$ paper-level slot 3 $\cdot$ primary & This reviewer halucinated the most and it spit out word salads. I am sure there were some original insights in there (the review was 5 times longer than the others), but they were clouded by verbosity \\[3pt]
\end{longtable}
}

\subsection*{W6: Trivial / nitpicking --- n = 16}

{\footnotesize
\begin{longtable}{@{}p{3.2cm}p{11.5cm}@{}}
\toprule
\textbf{Citation} & \textbf{Expert comment} \\
\midrule
\endfirsthead
\toprule
\textbf{Citation} & \textbf{Expert comment} \\
\midrule
\endhead
\bottomrule
\endlastfoot
P3 $\cdot$ GPT-5.2 $\cdot$ item 3 $\cdot$ primary & I believe it is just a yet another result the authors present which is essentially disconnected from the main narrative. \\[3pt]
P17 $\cdot$ GPT-5.2 $\cdot$ item 3 $\cdot$ primary & It seems in the end criticizing a typo in the caption of Figure 3, 1000 \textbackslash\{\}omega\textasciicircum{}\{-1\}\_\{ci\} --\textgreater{} 100 \textbackslash\{\}omega\textasciicircum{}\{-1\}\_\{ci\} but the comment overly criticizing. Just pointing out a typo is enough. \\[3pt]
P25 $\cdot$ Claude 4.5 $\cdot$ item 1 $\cdot$ primary & more of a technicallity. \\[3pt]
P26 $\cdot$ GPT-5.2 $\cdot$ item 4 $\cdot$ primary & seems like a technicallity to me. \\[3pt]
P42 $\cdot$ Gemini 3.0 $\cdot$ item 3 $\cdot$ primary & Pretty sure that sentence is just a typo, not a misunderstanding and validity risk. \\[3pt]
P46 $\cdot$ GPT-5.2 $\cdot$ item 2 $\cdot$ primary & Do not think this is particular relevant. Sure, there may be badly predicted values outside the range... \\[3pt]
P53 $\cdot$ Gemini 3.0 $\cdot$ item 2 $\cdot$ primary & It is ok observation but definitely does not limit the study and the tools. it just makes it specific. \\[3pt]
P59 $\cdot$ Claude 4.5 $\cdot$ item 2 $\cdot$ secondary & This pairs with Human reviewer 3 critique about acidity preservation however here the critique is more of a ``the authors should discuss this limitation'' point than a fatal flaw. \\[3pt]
P62 $\cdot$ Gemini 3.0 $\cdot$ item 1 $\cdot$ primary & The comment is valid but the paper presents this as a grand vision rather than something it achieves, that is why this comment is marginally signicant rather than significant. \\[3pt]
P66 $\cdot$ Claude 4.5 $\cdot$ item 1 $\cdot$ primary & When we fit data with a Lorentzian using MATLAB for example, MATLAB gives us the best fit and range of values which leads to, let's say, R2 value of 0.99. The point raised by the AI reviewer is minor with respect to the main message of the manuscript. \\[3pt]
P67 $\cdot$ Claude 4.5 $\cdot$ item 3 $\cdot$ primary & This is a very minor issue. In fact, sometimes I personally forget to include the statistical analysis of ML outcomes. \\[3pt]
P68 $\cdot$ Claude 4.5 $\cdot$ item 1 $\cdot$ primary & This is true, but unlikely to change any major conclusion \\[3pt]
P68 $\cdot$ Gemini 3.0 $\cdot$ item 5 $\cdot$ primary & These are indeed speculative assertions that can be easily dealt with. \\[3pt]
P72 $\cdot$ Claude 4.5 $\cdot$ paper-level slot 1 $\cdot$ primary & AI reviewer 1 noted that the GIAB ground-truth set contained only SVs within the ``high-confidence'' regions. I consider to be a relatively minor point that could be addressed by emphasizing that benchmarking results may not generalize to more difficult/repetitive parts of the genome. \\[3pt]
P72 $\cdot$ Gemini 3.0 $\cdot$ paper-level slot 2 $\cdot$ primary & An unused argument in the code, which is hepful but about as minor as a criticism can be. \\[3pt]
P73 $\cdot$ Gemini 3.0 $\cdot$ paper-level slot 1 $\cdot$ primary & AI Reviewer 1 was the only reviewer to point out that there are newer methods for detecting microsat mutations, but I felt that this was a minor point. \\[3pt]
\end{longtable}
}

\subsection*{W7: Technical term-of-art confusion --- n = 13}

{\footnotesize
\begin{longtable}{@{}p{3.2cm}p{11.5cm}@{}}
\toprule
\textbf{Citation} & \textbf{Expert comment} \\
\midrule
\endfirsthead
\toprule
\textbf{Citation} & \textbf{Expert comment} \\
\midrule
\endhead
\bottomrule
\endlastfoot
P25 $\cdot$ Gemini 3.0 $\cdot$ item 1 $\cdot$ primary & On first sight it seems like a good comment, but AI and TW are established as two different aspects of neuronal dynamics, the first is local, the second is spatially extended. \\[3pt]
P28 $\cdot$ Gemini 3.0 $\cdot$ item 4 $\cdot$ primary & There is a misunderstanding of what the term ''universal'' means here. ''Universality'' is mere a physics term that means the data collaps to a single curve when scaled. It does not mean it is universally valid. \\[3pt]
P42 $\cdot$ Gemini 3.0 $\cdot$ item 2 $\cdot$ primary & This seems a bit of a terminology/scope issue, not a validity issue. \\[3pt]
P54 $\cdot$ Gemini 3.0 $\cdot$ item 3 $\cdot$ primary & AI should be aware of the connotation that is accepted by the genomics community for the use of terminologies such is real time. \\[3pt]
P65 $\cdot$ Gemini 3.0 $\cdot$ item 2 $\cdot$ primary & INR might take a few seconds for training, but once the training is complete, then it is extremely fast during inference. \\[3pt]
P66 $\cdot$ Gemini 3.0 $\cdot$ item 3 $\cdot$ primary & Here the resolution depends on the width of the tip, not the excitation wavelength. \\[3pt]
P67 $\cdot$ Claude 4.5 $\cdot$ item 4 $\cdot$ primary & I think AI got this totally wrong. Since 6 photodetectors have transmission peaks at different at different wavelengths, they can provide a much richer training data which eventually leads to a more accurate edge detectioon. \\[3pt]
P67 $\cdot$ Gemini 3.0 $\cdot$ item 1 $\cdot$ primary & SCA acts like a monochromatic filter, MCA is more like an array filter. \\[3pt]
P67 $\cdot$ Gemini 3.0 $\cdot$ item 2 $\cdot$ primary & SCA output is more like the average of 6 photodiodes of the same type. MCA output is more like an output of 6 different type photodiodes, which can distinguish the edges more accurately. \\[3pt]
P72 $\cdot$ GPT-5.2 $\cdot$ item 1 $\cdot$ primary & This comment is in line with one of the human reviewer's criticisms about the lack of methodological detail included in the paper. I don't get the AI reviewer's complaint about requiring a reference genome though, given that DeBreak is an read alignment-based method for detecting and locating SVs. \\[3pt]
P77 $\cdot$ Claude 4.5 $\cdot$ item 2 $\cdot$ secondary & In most of the substrates, there are no stereogenic centers to consider. If the reviewer is referring to the newly formed stereogenic centers arising from the C--C bond cleavage (Table 3), this is a valid point; however, these centers are not related to the electrocyclization step. \\[3pt]
P77 $\cdot$ Gemini 3.0 $\cdot$ item 3 $\cdot$ secondary & While the terminology appears to be somewhat misleading, both C--H and C--C bond cleavages would not take place in the absence of the metal carbene. \\[3pt]
P2 $\cdot$ Gemini 3.0 $\cdot$ paper-level slot 3 $\cdot$ primary & The issue with the size-extensivity of the proposed approach is a very important one. However, it should not be an argument for rejecting this work. I believer, the authors should simply thoroughly discuss this limitation. Overall, I found all AI reviewer very good but quite critical and rather not differentiating issues stemming from quantum chemistry and machine learning. I think, the issues from the former should be discussed in the paper, but not be a major issue when evaluating the proposed machine learning approach. \\[3pt]
\end{longtable}
}

\subsection*{W8: Cites evidence from after the preprint --- n = 9}

{\footnotesize
\begin{longtable}{@{}p{3.2cm}p{11.5cm}@{}}
\toprule
\textbf{Citation} & \textbf{Expert comment} \\
\midrule
\endfirsthead
\toprule
\textbf{Citation} & \textbf{Expert comment} \\
\midrule
\endhead
\bottomrule
\endlastfoot
P1 $\cdot$ Claude 4.5 $\cdot$ item 1 $\cdot$ primary & While claims are actually fine, I marked it as not correct as the paper they use to support these claims appeared two years after the presented one. The field was evolving very fast and in 2022 the presented results could be considered as good. The only thing one could move on with is the claims of improving upon the state-of-the-art. The auhtors should clarify what exactly they mean with it. \\[3pt]
P5 $\cdot$ Claude 4.5 $\cdot$ item 1 $\cdot$ primary & Funnily, this AI reviewer references in the last quote the same paper published later in 2022. \\[3pt]
P20 $\cdot$ Claude 4.5 $\cdot$ item 1 $\cdot$ primary & This point is significant and valid, however it lacks context and chronology. While it is true that spectroscopic confirmation is crucial, photometric evidence can still be compelling, especially given the fact that this discovery is reported very early after JWST started operations. Given this context, reporting such a potentialy far-reaching photometric discovery is acceptable. Additionally, the AI reviewer provides evidence from papers from 2025 which is invalid. \\[3pt]
P20 $\cdot$ Claude 4.5 $\cdot$ item 2 $\cdot$ primary & Again, while this point is valid and significant, the evidence is based on knowledge obtained in the future related to when this preprint was reviewed. Although a human reviewer should have asked the authors to comment on the AGN contamination, the significance would have been less because at that time we didn't have the knowledge of this class of objects. \\[3pt]
P20 $\cdot$ Gemini 3.0 $\cdot$ item 1 $\cdot$ primary & Again chronology is not preserved and the cited evidence is after this preprint \\[3pt]
P21 $\cdot$ Claude 4.5 $\cdot$ item 1 $\cdot$ primary & The comments cite a paper published 2 years after this preprint which is not valid. \\[3pt]
P46 $\cdot$ Claude 4.5 $\cdot$ item 3 $\cdot$ primary & This needs to be seen in historical context. PolyBERT was very early. \\[3pt]
P20 $\cdot$ Claude 4.5 $\cdot$ paper-level slot 1 $\cdot$ primary & Yes, but the evidence was based on knowledge obtained several years after the publication of this preprint. \\[3pt]
P20 $\cdot$ Gemini 3.0 $\cdot$ paper-level slot 3 $\cdot$ primary & Yes, but the evidence was based on knowledge obtained several years after the publication of this preprint. \\[3pt]
\end{longtable}
}

\subsection*{W9: Over-inflates small code/text inconsistencies --- n = 9}

{\footnotesize
\begin{longtable}{@{}p{3.2cm}p{11.5cm}@{}}
\toprule
\textbf{Citation} & \textbf{Expert comment} \\
\midrule
\endfirsthead
\toprule
\textbf{Citation} & \textbf{Expert comment} \\
\midrule
\endhead
\bottomrule
\endlastfoot
P12 $\cdot$ Claude 4.5 $\cdot$ item 2 $\cdot$ secondary & The two Doppler estimates are not shown to be in serious conflict, because the variability estimate has no quoted uncertainty. \\[3pt]
P34 $\cdot$ Gemini 3.0 $\cdot$ item 2 $\cdot$ primary & This one could be an unit error in the code, over-interpreted by AI-reviewer. \\[3pt]
P34 $\cdot$ Gemini 3.0 $\cdot$ item 4 $\cdot$ primary & This could be also the unit error in the code. \\[3pt]
P42 $\cdot$ GPT-5.2 $\cdot$ item 2 $\cdot$ primary & I think the AI is overstating the implications of small deviations between code and text that occur over the lifetime of a scientific paper. \\[3pt]
P42 $\cdot$ GPT-5.2 $\cdot$ item 5 $\cdot$ primary & Same comment as item 2. \\[3pt]
P46 $\cdot$ GPT-5.2 $\cdot$ item 4 $\cdot$ primary & This is most likely not a real problem and solved upon final submission. \\[3pt]
P54 $\cdot$ Claude 4.5 $\cdot$ item 5 $\cdot$ primary & Quote: ''Functional annotation of the assembled plasmid genomes using BugSeq revealed that KPC-2 and KPC-14 were both located on IncN plasmids, which showed 99.7\% identity according to sequence alignments''-overstatement by AI and overanalysis \\[3pt]
P59 $\cdot$ Claude 4.5 $\cdot$ item 3 $\cdot$ secondary & This feels very much AI generated critique that identifies a real asymmetry in methodology but fails to evalute whetjer it matters qunatitatively. the cement sector's contribution to total TPAE is so small that even substantial provincial variation in cement toxicity would not meaningfully alter the conclusions. \\[3pt]
P59 $\cdot$ GPT-5.2 $\cdot$ item 1 $\cdot$ secondary & The within-sector variation is modest compared to the \textasciitilde{}5 fold between-sector differences. So while the two aggregation methods would yield different absolute TU values, the sector ranking is very unlikely to change. However, the reviewer is correct that the authors don't demonstrate this, a sensitivity analysis would be needed to confirm. \\[3pt]
\end{longtable}
}

\subsection*{W10: Criticizes what authors already flagged as a limitation --- n = 6}

{\footnotesize
\begin{longtable}{@{}p{3.2cm}p{11.5cm}@{}}
\toprule
\textbf{Citation} & \textbf{Expert comment} \\
\midrule
\endfirsthead
\toprule
\textbf{Citation} & \textbf{Expert comment} \\
\midrule
\endhead
\bottomrule
\endlastfoot
P3 $\cdot$ Claude 4.5 $\cdot$ item 2 $\cdot$ primary & The authors state explicitly that it is a proof-of-concept study. Transferability claims are somewhat unrelated to the main narrative but it does not make the reviewer's claim correct. \\[3pt]
P46 $\cdot$ GPT-5.2 $\cdot$ item 1 $\cdot$ primary & It is a problem that is also already mentioned by the authors. It is up to the reviewer to judge if that invalidates everything. I would not think so. \\[3pt]
P60 $\cdot$ Gemini 3.0 $\cdot$ item 2 $\cdot$ secondary & The core criticism is valid but should be calibrated against what the paper actually discusses later. The deeper methodological question whether to fit separate models for the COVID period is something the authors themselves flag as future work, constrained by insufficient data (only 2 years of COVID-period observations for regional-level estimation). \\[3pt]
P60 $\cdot$ Gemini 3.0 $\cdot$ item 3 $\cdot$ secondary & This is a well-articulated critique, but it partly overlaps with Item 2 and with what the authors themselves discuss. The paper actually acknowledges this directly in the text quoted by the reviewer and in the discussion around Figure 4, where the authors present evidence suggesting the two hazards do interact and explicitly state they assumed independence. \\[3pt]
P62 $\cdot$ Gemini 3.0 $\cdot$ item 3 $\cdot$ secondary & The algorithm is essentially NMF and authors stated that they used the existing NMF package already available. The way to use NMF in fiber photometry context is new, so I think this comment, which is regarding code implementation, is marginal. \\[3pt]
P62 $\cdot$ Gemini 3.0 $\cdot$ item 4 $\cdot$ secondary & The key point of this study is not in NMF, nor the authors does not claim NMF as their developed method. The way to use NMF in fiber photometry context is new, so I think this comment is marginal. \\[3pt]
\end{longtable}
}

\subsection*{W11: AI misreads a figure or caption --- n = 6}

{\footnotesize
\begin{longtable}{@{}p{3.2cm}p{11.5cm}@{}}
\toprule
\textbf{Citation} & \textbf{Expert comment} \\
\midrule
\endfirsthead
\toprule
\textbf{Citation} & \textbf{Expert comment} \\
\midrule
\endhead
\bottomrule
\endlastfoot
P3 $\cdot$ Gemini 3.0 $\cdot$ item 1 $\cdot$ primary & For example, Figure S10 does not show barriers of 40-50 kJ/mol. Furthermore, Figure S13 shows every barriers of a comparable height to previous studies. Finally, the chosen DFT may be a valid critique but the results should be interpreted within this setting if it is not dramatically wrong. \\[3pt]
P19 $\cdot$ Gemini 3.0 $\cdot$ item 3 $\cdot$ primary & Maybe neglecting Figure 4. \\[3pt]
P35 $\cdot$ Claude 4.5 $\cdot$ item 3 $\cdot$ primary & review misidentifies the optimal threshold as t=0.4 when Figure 5c shows the peak at t=0.3. \\[3pt]
P35 $\cdot$ Claude 4.5 $\cdot$ item 4 $\cdot$ primary & error bars are present in figure2c (small gray lines on each bar) \\[3pt]
P66 $\cdot$ GPT-5.2 $\cdot$ item 4 $\cdot$ primary & I do not agree with the AI reviewer again. The uniformity in Fig3. d and e removes the questions about the defects. Since these were mechanical exfoliated samples, strain criticism is also not relevant. \\[3pt]
P77 $\cdot$ Claude 4.5 $\cdot$ item 3 $\cdot$ secondary & The reviewer may have misunderstood Figure 2 in the main text. \\[3pt]
\end{longtable}
}

\subsection*{W12: AI misquotes or fabricates a verbatim quote --- n = 5}

{\footnotesize
\begin{longtable}{@{}p{3.2cm}p{11.5cm}@{}}
\toprule
\textbf{Citation} & \textbf{Expert comment} \\
\midrule
\endfirsthead
\toprule
\textbf{Citation} & \textbf{Expert comment} \\
\midrule
\endhead
\bottomrule
\endlastfoot
P29 $\cdot$ GPT-5.2 $\cdot$ item 1 $\cdot$ primary & I doubt AI run the model. \\[3pt]
P49 $\cdot$ Gemini 3.0 $\cdot$ item 1 $\cdot$ secondary & Instead of quoting the author's exact words, this AI reviewer incorrectly quoted the author and then claimed that the author's claims were incorrect. It's bit surprising that AI reviewers can make such mistakes, and I think we need to be aware that they can make such simple mistakes. \\[3pt]
P65 $\cdot$ Gemini 3.0 $\cdot$ item 4 $\cdot$ primary & Again, AI made a mistake by comparing the sentences taken from different parts of the manuscript talking about different parts of the process. \\[3pt]
P77 $\cdot$ GPT-5.2 $\cdot$ item 3 $\cdot$ secondary & I could not find CuTe in both main text and supplementary information. \\[3pt]
P77 $\cdot$ Gemini 3.0 $\cdot$ item 1 $\cdot$ secondary & I could not observe the discrepancies pointed out by the reviewer. \\[3pt]
\end{longtable}
}

\subsection*{W13: AI misses supplementary content --- n = 3}

{\footnotesize
\begin{longtable}{@{}p{3.2cm}p{11.5cm}@{}}
\toprule
\textbf{Citation} & \textbf{Expert comment} \\
\midrule
\endfirsthead
\toprule
\textbf{Citation} & \textbf{Expert comment} \\
\midrule
\endhead
\bottomrule
\endlastfoot
P6 $\cdot$ Gemini 3.0 $\cdot$ item 1 $\cdot$ secondary & very detailed in the SI: The RF and DB 0 ?eld maps were acquired for each subject for the pulse design and shimming using 2D interferometric turbo-FLASH [30] at 5 mm isotropic resolution and 3D multi-echo (TE = 1.6/3.5/6 ms) GRE acquisitions at 2.5 mm isotropic resolution respectively. Sequence parameters for the T 2 -weighted variable ?ip angle turbo spin-echo acquisition were: resolution = 0.55 $\times$ 0.55 $\times$ 0.55 mm 3 , TR = 6 s, TA = 13 min, GRAPPA = 3 $\times$ 2, TE = 301 ms, matrix = 400 $\times$ 400 $\times$ 320, bandwidth = 250 Hz/pixel. \\[3pt]
P37 $\cdot$ GPT-5.2 $\cdot$ item 1 $\cdot$ secondary & The first comment is correct. The second comment is correct based on what is written in the main text. But if you refer to the Supplementary Figure 3, then we can see the number 281 in the main text is a typo. The actual number is 69281. It seems the AI Reviewer did not check the supplementary material. \\[3pt]
P53 $\cdot$ Claude 4.5 $\cdot$ item 1 $\cdot$ primary & In supplementary \\[3pt]
\end{longtable}
}

\subsection*{W14: Ignores authors' own cited prior work --- n = 2}

{\footnotesize
\begin{longtable}{@{}p{3.2cm}p{11.5cm}@{}}
\toprule
\textbf{Citation} & \textbf{Expert comment} \\
\midrule
\endfirsthead
\toprule
\textbf{Citation} & \textbf{Expert comment} \\
\midrule
\endhead
\bottomrule
\endlastfoot
P28 $\cdot$ Claude 4.5 $\cdot$ item 2 $\cdot$ primary & The authors make the claim about correlation of nutrient concentration and interaction strength based on their previous work. This critique is similar to the one of the Human Reviewers. \\[3pt]
P28 $\cdot$ Gemini 3.0 $\cdot$ item 2 $\cdot$ primary & The authors make this claim based on their previous work. Similar to one of the human reviewers, the AI reviewer did not read those cited papers. \\[3pt]
\end{longtable}
}

\subsection*{W15: AI misreads a table --- n = 1}

{\footnotesize
\begin{longtable}{@{}p{3.2cm}p{11.5cm}@{}}
\toprule
\textbf{Citation} & \textbf{Expert comment} \\
\midrule
\endfirsthead
\toprule
\textbf{Citation} & \textbf{Expert comment} \\
\midrule
\endhead
\bottomrule
\endlastfoot
P32 $\cdot$ Claude 4.5 $\cdot$ item 2 $\cdot$ primary & The reviewer's assessment seems to be based on an over-interpretation of Table 1, where a typographical error in the values may have led to an incorrect conclusion. \\[3pt]
\end{longtable}
}

\subsection*{W16: Cannot analyze figures --- only text --- n = 1}

{\footnotesize
\begin{longtable}{@{}p{3.2cm}p{11.5cm}@{}}
\toprule
\textbf{Citation} & \textbf{Expert comment} \\
\midrule
\endfirsthead
\toprule
\textbf{Citation} & \textbf{Expert comment} \\
\midrule
\endhead
\bottomrule
\endlastfoot
P6 $\cdot$ Gemini 3.0 $\cdot$ item 5 $\cdot$ primary & This could be argued better, but to do this the AI would need to be able to analyze the images. It went by text only, did not download the figures from Figshare, so it gave boilerplate feedback. \\[3pt]
\end{longtable}
}

\subsection*{W\_unspecified: Residual: AI judged Not Correct without specific reason --- n = 10}

{\footnotesize
\begin{longtable}{@{}p{3.2cm}p{11.5cm}@{}}
\toprule
\textbf{Citation} & \textbf{Expert comment} \\
\midrule
\endfirsthead
\toprule
\textbf{Citation} & \textbf{Expert comment} \\
\midrule
\endhead
\bottomrule
\endlastfoot
P28 $\cdot$ Gemini 3.0 $\cdot$ item 5 $\cdot$ primary & I don't think this is correct. \\[3pt]
P29 $\cdot$ Gemini 3.0 $\cdot$ item 1 $\cdot$ primary & Complete fail here. \\[3pt]
P30 $\cdot$ GPT-5.2 $\cdot$ item 4 $\cdot$ primary & I am not sure where this code is. \\[3pt]
P34 $\cdot$ Gemini 3.0 $\cdot$ item 1 $\cdot$ secondary & I did not verify this in detail by running the Python code, and I assume the authors did not make this mistake. However, if this is indeed the case, it would be a significant issue. \\[3pt]
P3 $\cdot$ Claude 4.5 $\cdot$ paper-level slot 2 $\cdot$ primary & I found this reviewer rather bad. \\[3pt]
P3 $\cdot$ Gemini 3.0 $\cdot$ paper-level slot 3 $\cdot$ primary & I found this reviewer rather bad. \\[3pt]
P9 $\cdot$ GPT-5.2 $\cdot$ paper-level slot 1 $\cdot$ primary & Human review 2 was not provided in the checklist. For the last question here I added AI reviewer 3 however it's a bit borderline, I think the 3 AI reviews were not very good. \\[3pt]
P77 $\cdot$ Gemini 3.0 $\cdot$ paper-level slot 1 $\cdot$ primary & AI Reviewer 1 offered the weakest assessment among all reviewers. \\[3pt]
P79 $\cdot$ GPT-5.2 $\cdot$ paper-level slot 1 $\cdot$ primary & AI reviewer 1 does not seem to have a fundamental understanding of the key discovery of the manuscript. \\[3pt]
P80 $\cdot$ GPT-5.2 $\cdot$ paper-level slot 1 $\cdot$ primary & it misses the main point of the paper. \\[3pt]
\end{longtable}
}

\subsection*{S1: Statistical / methodology rigor --- n = 45}

{\footnotesize
\begin{longtable}{@{}p{3.2cm}p{11.5cm}@{}}
\toprule
\textbf{Citation} & \textbf{Expert comment} \\
\midrule
\endfirsthead
\toprule
\textbf{Citation} & \textbf{Expert comment} \\
\midrule
\endhead
\bottomrule
\endlastfoot
P7 $\cdot$ Gemini 3.0 $\cdot$ item 2 $\cdot$ primary & Here in general the problem is that the fit is of bad quality \\[3pt]
P35 $\cdot$ GPT-5.2 $\cdot$ item 1 $\cdot$ primary & The PPI trainer has no validation split (only train/test), uses test metrics for model selection (misleadingly named best\_valid\_f1) \\[3pt]
P60 $\cdot$ Claude 4.5 $\cdot$ item 3 $\cdot$ secondary & Even if the overall sample size were adequate, treating correlated regional observations within the same event as independent violates the K-S test assumptions. This is a genuinely useful statistical critique that neither the first reviewer in this set nor the human reviewers identified. \\[3pt]
P60 $\cdot$ Gemini 3.0 $\cdot$ item 1 $\cdot$ secondary & This is an outstanding technical critique. It's worth noting that the attributable AF itself is independent of the baseline, so the relative risk estimates are unaffected but the absolute death counts, which are the paper's primary output, scale directly with deathexp. \\[3pt]
P60 $\cdot$ Gemini 3.0 $\cdot$ item 5 $\cdot$ secondary & The 1.7--3.8x fold increases for heatwaves and 2.1--3.5x for cold snaps are presented as key results demonstrating the severity of compound crises, but as the reviewer notes, these numbers are mathematically inevitable once COVID-19 deaths exist. This connects to Item 3's point about the paper not truly measuring compound/interactive effects. \\[3pt]
P61 $\cdot$ Claude 4.5 $\cdot$ item 3 $\cdot$ secondary & this is logically paired with Item 2 and together they form a coherent critique: the provincial methodology is underspecified (Item 2) and unvalidated (Item 3). It's possible that the Pearson correlation in Figure S3 does include some provincial-level comparisons but again since its not explicit, authors needs to work on the presentation. \\[3pt]
P61 $\cdot$ Claude 4.5 $\cdot$ item 4 $\cdot$ secondary & The authors should clarify: did they combine bounds simultaneously (which would overestimate uncertainty for independent sources), use Monte Carlo sampling, or apply some other approach? This is a fair ask for methodological transparency even if the practical impact on conclusions is small. \\[3pt]
P61 $\cdot$ GPT-5.2 $\cdot$ item 2 $\cdot$ secondary & The paper would benefit from a supplementary table of fitted parameters, a description of the ARIMA specification procedure, and explicit statement of the TMREL distribution used. \\[3pt]
P61 $\cdot$ GPT-5.2 $\cdot$ item 3 $\cdot$ secondary & The reviewer correctly distinguishes between association (correlation) and agreement (accuracy in magnitude), citing the Bland-Altman framework which is the standard methodological reference for exactly this issue \\[3pt]
P61 $\cdot$ GPT-5.2 $\cdot$ item 4 $\cdot$ secondary & This critique pairs well with Review 1's Item 4 (incomplete uncertainty propagation) but is far more precise. The critique doesn't necessarily invalidate the central qualitative findings (aging as dominant driver), but it does undermine the quantitative confidence statements that the paper emphasizes throughout. \\[3pt]
P63 $\cdot$ Gemini 3.0 $\cdot$ item 4 $\cdot$ primary & The power criticism is very correct and very significant. On the time, comments are correct but the mentioned challenges can be overcome. Howecer the enegy part is the fundamental citicism that should be enough to reject the paper. \\[3pt]
P65 $\cdot$ GPT-5.2 $\cdot$ item 4 $\cdot$ primary & I agree with the AI: ''a great deal'' does not mean anything from a scientific point of view. They should have quantified their finding. \\[3pt]
P68 $\cdot$ Gemini 3.0 $\cdot$ item 2 $\cdot$ primary & Given that there are categorical PLS options, this is a a reasonable criticism. \\[3pt]
P73 $\cdot$ GPT-5.2 $\cdot$ item 1 $\cdot$ primary & These methods do appear to be missing, and could be consequential. \\[3pt]
P74 $\cdot$ Claude 4.5 $\cdot$ item 2 $\cdot$ primary & I agree that data leakage is possible if the same regions are included in both training and test sets. A safer approach would be to split the genome up into segments that are partitioned into the training and test sets. \\[3pt]
P3 $\cdot$ GPT-5.2 $\cdot$ paper-level slot 1 $\cdot$ primary & This AI reviewer recognized ambiguity in the description of methods used for certain calculations and asks for the clarification for the computed coverages. \\[3pt]
P4 $\cdot$ Gemini 3.0 $\cdot$ paper-level slot 2 $\cdot$ primary & Essentially, I find again great that also this reviewer questions if random splitting is a good strategy, similar to the first AI reviewer. \\[3pt]
P5 $\cdot$ Gemini 3.0 $\cdot$ paper-level slot 1 $\cdot$ primary & I found this AI reviewer quite good. They were very constructive and concise. Their arguments are close to those I would make. Similar to the third reviewer, the most significant contribution, in my opinion, was pointing out the lack of a proper baseline comparison and the absence of variance estimates in the computed errors of the predicted values. \\[3pt]
P14 $\cdot$ Claude 4.5 $\cdot$ paper-level slot 1 $\cdot$ secondary & Yes, Potential confounding of circadian rhythm detection by ongoing therapeutic stimulation. Human reviewer only consider the relation with the data analysis issues. \\[3pt]
P15 $\cdot$ Claude 4.5 $\cdot$ paper-level slot 2 $\cdot$ primary & Absence of multiple comparisons \\[3pt]
P15 $\cdot$ Claude 4.5 $\cdot$ paper-level slot 2 $\cdot$ secondary & [English translation from Korean] Argued the problem of the absence of a healthy control group in the most systematic way. \\[3pt]
P15 $\cdot$ GPT-5.2 $\cdot$ paper-level slot 1 $\cdot$ primary & Yes, data imputation methodological problem and inconsistencies; multiple comparisons absence. \\[3pt]
P15 $\cdot$ Gemini 3.0 $\cdot$ paper-level slot 3 $\cdot$ primary & Lack of analyses on EEG microstructures; data imputation problem \\[3pt]
P16 $\cdot$ Claude 4.5 $\cdot$ paper-level slot 3 $\cdot$ primary & granger causality group level significance but not individual level \\[3pt]
P24 $\cdot$ GPT-5.2 $\cdot$ paper-level slot 3 $\cdot$ primary & catching the normative atlas issue missed by the humans. nice catching the lack of spurious correlations null model. \\[3pt]
P35 $\cdot$ Claude 4.5 $\cdot$ paper-level slot 3 $\cdot$ primary & The utilization rate metric mathematical error and the suspicious standard deviation entry in Table S2 were unique observations. \\[3pt]
P35 $\cdot$ GPT-5.2 $\cdot$ paper-level slot 2 $\cdot$ primary & The evaluation protocol violation (missing validation split, model selection on test set) provides a complementary angle to AI Review 1's data leakage finding. The BioSNAP benchmark limitation and code reproducibility issues were unique to AI Reviewer 2. \\[3pt]
P37 $\cdot$ GPT-5.2 $\cdot$ paper-level slot 3 $\cdot$ secondary & It was the only reviewer who noticed: - Class imbalance in the dataset affects AUPRC, which can undermine the SOTA claim of the paper. - The submitted code does not accurately calculate AUROC \\[3pt]
P42 $\cdot$ Gemini 3.0 $\cdot$ paper-level slot 2 $\cdot$ primary & AI reviewer 2 found the collinearity of habitat fragmentation and amount and had a good suggestion for sensitivity analysis. \\[3pt]
P47 $\cdot$ Claude 4.5 $\cdot$ paper-level slot 3 $\cdot$ primary & AI \#3 uniquely pointed out that the entropy equation is mathematically incorrect. \\[3pt]
P47 $\cdot$ Gemini 3.0 $\cdot$ paper-level slot 1 $\cdot$ primary & AI \#1 was the only reviewer who raised a critical question regarding the intra-hamming distance, and argued that its magnitude is too large to support stable cryptographic key generation. \\[3pt]
P49 $\cdot$ Claude 4.5 $\cdot$ paper-level slot 1 $\cdot$ secondary & AI reviewer 1 pointed out a statistical flaw in the distance-dependent cross-validation method. Specifically, test regions adjacent to the boundary of the 75\% training set and the 25\% test set still share spatial autocorrelation, which can potentially inflate the model's performance estimates. \\[3pt]
P50 $\cdot$ GPT-5.2 $\cdot$ paper-level slot 1 $\cdot$ primary & Yes, AI reveiwer 1 found out about SNR which all other reviewers failed to catch. \\[3pt]
P50 $\cdot$ GPT-5.2 $\cdot$ paper-level slot 1 $\cdot$ secondary & AI Reviewer 1's review included details such as the SNR that was not in other reviews, human or otherwise. \\[3pt]
P52 $\cdot$ GPT-5.2 $\cdot$ paper-level slot 3 $\cdot$ primary & Yes. Issue of experimental setting on ``control''. \\[3pt]
P52 $\cdot$ GPT-5.2 $\cdot$ paper-level slot 3 $\cdot$ secondary & confounding factors related to contralateral effects of the drug that was used unilaterally, with the contralateral eye as a control. Also, this reviewer mentioned the code that was used in the paper. \\[3pt]
P55 $\cdot$ Claude 4.5 $\cdot$ paper-level slot 3 $\cdot$ primary & AI reviewer 3 overestimated the age of the population and had similar comments as the other AI reviewers. it did cough more statistics and model flaws than all the human comments. same applies for AI 1 and AI2. \\[3pt]
P55 $\cdot$ Gemini 3.0 $\cdot$ paper-level slot 1 $\cdot$ primary & MS methodology in for example catching the imputation in the data analysis. this was not cought by the human reviewer. \\[3pt]
P67 $\cdot$ Claude 4.5 $\cdot$ paper-level slot 3 $\cdot$ primary & AIR 3 caught the lack of statistical analysis, which was not caught by other reviewers. AIR2 and AIR3 did a very good job in terms finding the different resolution used in SCA and MCA, which I hope was fixed later by the authors. \\[3pt]
P68 $\cdot$ Gemini 3.0 $\cdot$ paper-level slot 2 $\cdot$ primary & The ordinal data issue was raised by AI reviewer 2 and not by others. \\[3pt]
P79 $\cdot$ Claude 4.5 $\cdot$ paper-level slot 2 $\cdot$ primary & AI reviewers 2 and 3 both made critical assessment regarding the causality \\[3pt]
P79 $\cdot$ Gemini 3.0 $\cdot$ paper-level slot 3 $\cdot$ primary & AI reviewers 2 and 3 both made critical assessment regarding the causality. AI Reviewer 3 was more in depth than AI reviewer 2. \\[3pt]
P80 $\cdot$ Claude 4.5 $\cdot$ paper-level slot 3 $\cdot$ secondary & Yes. AI Reviewer 3 gave reasonable comments on the simulation setting, theoretical assumption, and uncertainty in real data analysis that other reviewers did not catch. \\[3pt]
P85 $\cdot$ GPT-5.2 $\cdot$ paper-level slot 1 $\cdot$ primary & Library preparation batch effects \\[3pt]
P85 $\cdot$ Gemini 3.0 $\cdot$ paper-level slot 3 $\cdot$ primary & AI Reviewer 3 was very focused on data analysis (computational) procedures and almost all comments were distinct from the other reviewers (human and AI). \\[3pt]
\end{longtable}
}

\subsection*{S2: Code reading --- n = 28}

{\footnotesize
\begin{longtable}{@{}p{3.2cm}p{11.5cm}@{}}
\toprule
\textbf{Citation} & \textbf{Expert comment} \\
\midrule
\endfirsthead
\toprule
\textbf{Citation} & \textbf{Expert comment} \\
\midrule
\endhead
\bottomrule
\endlastfoot
P1 $\cdot$ GPT-5.2 $\cdot$ item 4 $\cdot$ primary & I find it quite impressive that this reviewer actually looks at the code to understand whether some parts of the paper are not properly discussed. I am personally not doing it on a regular basis as it would consume a substantial amount of time. \\[3pt]
P6 $\cdot$ Gemini 3.0 $\cdot$ item 4 $\cdot$ primary & Impressive that the AI looked at the code and realized there is nothing about the MRI reconstruction. It missed that some of the links are broken even for the ANT-R task. \\[3pt]
P35 $\cdot$ Gemini 3.0 $\cdot$ item 1 $\cdot$ primary & Code verification fully confirms the data leakage claim. The DPI link trainer (trainer\_dpi.py) also has a bug using test metrics for model selection despite having a validation set. \\[3pt]
P37 $\cdot$ Claude 4.5 $\cdot$ item 3 $\cdot$ secondary & I checked the validity of this item using Claude Code. \\[3pt]
P37 $\cdot$ GPT-5.2 $\cdot$ item 3 $\cdot$ secondary & I checked the validity of this item using Claude Code. \\[3pt]
P37 $\cdot$ GPT-5.2 $\cdot$ item 4 $\cdot$ secondary & I checked the validity of this item using Claude Code. \\[3pt]
P37 $\cdot$ Gemini 3.0 $\cdot$ item 2 $\cdot$ secondary & I didn't manually go through the code to verify the claim. Instead, I asked Claude Opus 4.6 to verify the claim and it said the claim was correct. \\[3pt]
P49 $\cdot$ GPT-5.2 $\cdot$ item 2 $\cdot$ secondary & Unlike human reviewers, this reviewer actually visited the GitHub open source code repository the authors listed in their papers, dissected and analyzed the Python scripts to identify problems in the analysis. \\[3pt]
P50 $\cdot$ Gemini 3.0 $\cdot$ item 2 $\cdot$ primary & Human reviewer would not open up the raw code to find this out - if this code is correct, basically the core data reported in this paper is wrong. Authors in the code said ``It appears that delay is needed in order not to clog the port'' meaning it's not possible to wirelessly transmit high frequency information. If this is true, Figure 4, for example, would mean the authors measured in wired manner and falsely claimed that their system is wireless (which is one key aspect of this paper). \\[3pt]
P67 $\cdot$ Claude 4.5 $\cdot$ item 1 $\cdot$ primary & Same as the AI reviewer-2, this reviewer also caught the same resolution issue. Maybe this is one of the a few advantages over the human reviewers. AI is good at reviewing the codes and finding bugs. \\[3pt]
P67 $\cdot$ GPT-5.2 $\cdot$ item 2 $\cdot$ primary & I didn't examine their code, but it looks like it wasn't a fair comparison. I hope the authors later fixed this. \\[3pt]
P74 $\cdot$ Gemini 3.0 $\cdot$ item 1 $\cdot$ primary & The AI reviewer seems to be correct here, although I did not go through the entireity of the code so I cannot be certain that this check was not done elsewhere. \\[3pt]
P1 $\cdot$ GPT-5.2 $\cdot$ paper-level slot 3 $\cdot$ primary & I think this reviewer also more clearly identified the systematic issue in the evaluation of the finite-size correction. I was also impressed by how this reviewer examined the code, identifying inconsistencies in the mathematical formulations. \\[3pt]
P15 $\cdot$ GPT-5.2 $\cdot$ paper-level slot 1 $\cdot$ secondary & [English translation from Korean] Code-level analysis is unique. These are problems that cannot be found by reading the paper text alone -- they are the result of analyzing the actual codebase. \\[3pt]
P23 $\cdot$ Claude 4.5 $\cdot$ paper-level slot 3 $\cdot$ primary & Good job in checking the code and finding the discrepancies. \\[3pt]
P23 $\cdot$ GPT-5.2 $\cdot$ paper-level slot 2 $\cdot$ primary & Good job in checking the code and finding the discrepancies. \\[3pt]
P24 $\cdot$ Gemini 3.0 $\cdot$ paper-level slot 2 $\cdot$ primary & nice checking the code and catching the normative atlas issue missed by the humans. \\[3pt]
P25 $\cdot$ GPT-5.2 $\cdot$ paper-level slot 3 $\cdot$ primary & Thoroughly checking the code and whether it can reproduce the results. \\[3pt]
P26 $\cdot$ GPT-5.2 $\cdot$ paper-level slot 2 $\cdot$ primary & good job on noticing and calling up the circularity in the work. Alsoo very useful for checking the code. \\[3pt]
P27 $\cdot$ Claude 4.5 $\cdot$ paper-level slot 2 $\cdot$ primary & human reviewers, especially 1, were very good. AI only exceeded them in checking the code. \\[3pt]
P27 $\cdot$ GPT-5.2 $\cdot$ paper-level slot 1 $\cdot$ primary & human reviewers, especially 1, were very good. AI only exceeded them in checking the code. \\[3pt]
P27 $\cdot$ Gemini 3.0 $\cdot$ paper-level slot 3 $\cdot$ primary & human reviewers, especially 1, were very good. AI only exceeded them in checking the code. \\[3pt]
P35 $\cdot$ Gemini 3.0 $\cdot$ paper-level slot 1 $\cdot$ primary & The graph structure data leakage mechanism and the swapped EM step naming in code were unique. The co-training critique overlaps with others but adds unique code evidence. \\[3pt]
P49 $\cdot$ GPT-5.2 $\cdot$ paper-level slot 2 $\cdot$ secondary & AI Reviewer 2 went beyond merely reading the text of the manuscript and providing logical critiques. It delved directly into the actual Python code scripts, data arrays, and references of the original datasets provided by the authors, accurately pinpointing critical discrepancies between what was claimed in the paper (methods) and what was actually executed (code or data). - False description of the cross-validation code: The manuscript claims to have performed 1,000 random iterations for distance-dependent cross-validation, but AI Reviewer 2 analyzed the provided Python code and revealed that it actually executed only 68 deterministic splits corresponding to the 68 brain regions. - Mismatch in the number of predictors: The methodology explicitly states that nine global connectome predictors were used, but AI Reviewer 2 found that only seven predictors were actually enumerated in the text and implemented in the code. - Age range error in ENIGMA data: Contrary to the authors' assertion that all cortical abnormality maps were collected from ''adult patients,'' the cited ENIGMA ASD dataset actually comprises participants aged 2 to 64 years, highlighting that unaddressed developmental lifespan effects could heavily distort the results. \\[3pt]
P50 $\cdot$ Gemini 3.0 $\cdot$ paper-level slot 2 $\cdot$ primary & Yes, AI Reviewer 2 found a problem that can be raised from raw code - which if true would flip all the claims this paper made on their capability of wireless system. \\[3pt]
P60 $\cdot$ Gemini 3.0 $\cdot$ paper-level slot 1 $\cdot$ secondary & This reviewer went directly into the code to uncover a concrete computational bias that no other reviewer could have found from reading the manuscript alone. It also identified opposing directions of bias, harvesting potentially overestimating temperature deaths while pandemic context underestimates them. It questioned that compound term was used only in wordings but not analytically. \\[3pt]
P67 $\cdot$ Gemini 3.0 $\cdot$ paper-level slot 1 $\cdot$ secondary & I find the review of AI reviewer 1 to be the best among all the human and AI reviews. One thing I would like to mention when comparing to human reviews is that this AI looked thoroughly into the code to find inconsistencies. \\[3pt]
P74 $\cdot$ Gemini 3.0 $\cdot$ paper-level slot 3 $\cdot$ primary & AI Reviewer 3 dug into the code and noticed the possiblity that ``hidden split reads'' are not actually being detected in the same way as described in the paper. \\[3pt]
\end{longtable}
}

\subsection*{S3: Specialized niche field catch --- n = 27}

{\footnotesize
\begin{longtable}{@{}p{3.2cm}p{11.5cm}@{}}
\toprule
\textbf{Citation} & \textbf{Expert comment} \\
\midrule
\endfirsthead
\toprule
\textbf{Citation} & \textbf{Expert comment} \\
\midrule
\endhead
\bottomrule
\endlastfoot
P61 $\cdot$ GPT-5.2 $\cdot$ item 1 $\cdot$ secondary & This critique substantially improves upon Human Reviewer 1's Item 20 (which asked about the formula's basis).If nitrate data were unavailable from the CMIP6 models used, the authors should have stated this as a limitation and discussed its implications for their estimates. \\[3pt]
P64 $\cdot$ GPT-5.2 $\cdot$ item 4 $\cdot$ primary & this is a good catch. retrieving the complex field (with the imaginary part) is very ill-posed and to the best of my knowledge there is no method to do that well with multimode fiber propagation. arbitrary optical fields is for sure not within the scope of the paper and the method has nothing to do with the complex field. the authors state the relation as amplitude-to-amplitude, so this part is accurate but for sure not enough to extrapolate this huge leap. \\[3pt]
P66 $\cdot$ GPT-5.2 $\cdot$ item 3 $\cdot$ primary & It is true that the bright-versus-dark exciton interpretation is not adequately supported and we do not know the exact illumination angle. But the authors claim a bright exciton dominance, so it should be close to a normal excitation. We should remember that the rapid spatial variation of the field at the tip apex is also required for dark excitons. If they were moving their tip slowly, then why would they have a serious dark excitons?? \\[3pt]
P66 $\cdot$ Gemini 3.0 $\cdot$ item 1 $\cdot$ primary & I agree with the AI reviewer that the PDM with their implementation is not the most accurate (they should have used the layered medium Green's functions and take care of the inhomogeneous background). However, this doesn't mean that the main findings of these researchers were wrong. A more accurate numerical model would generate results with smaller error limits, that's it. \\[3pt]
P73 $\cdot$ GPT-5.2 $\cdot$ item 3 $\cdot$ primary & I don't quite follow the argument about the GWAS results here, but the broader point about homopolymer errors is valid. \\[3pt]
P77 $\cdot$ GPT-5.2 $\cdot$ item 2 $\cdot$ primary & Since the substraste used for the synthesis of compound 4v contains stereogenic center, the issue of diastereoselectivity is important. Presumaby, mixture of diastereomers were formed in compound 4v. It was not mentioned by the human reviewers. \\[3pt]
P77 $\cdot$ Gemini 3.0 $\cdot$ item 2 $\cdot$ secondary & Although the formation of a seven-membered ring via an 8$\pi$ electrocyclization is relatively rare, such transformations have been reported. Nevertheless, I agree with the reviewer that the manuscript would be more impactful if additional experimental or computational evidence were provided to support the proposed mechanism. \\[3pt]
P1 $\cdot$ Claude 4.5 $\cdot$ paper-level slot 2 $\cdot$ primary & I think this reviewer more clearly identified the systematic issue in the evaluation of the finite-size correction. \\[3pt]
P2 $\cdot$ Claude 4.5 $\cdot$ paper-level slot 1 $\cdot$ primary & I found AI reviewer quite good. This identified important issues such as the lack of statistical sampling or the use of the minimal basis set, which may weaken the presented results. Also pointing out to a rather unfair comparison to GAP potentials appears to be raised by this AI reviewer but not the human ones, which I found very good. Overall, I found every single point raised by this AI reviewer as important and not necessarily found them in the human reviewer. \\[3pt]
P2 $\cdot$ GPT-5.2 $\cdot$ paper-level slot 2 $\cdot$ primary & This reviewer raises an issue concerning how close the provided results are really to the reference on the Zundel-cation benchmark. I found it quite important for the scope of this paper. And it appears not to be raised by the human reviewer. Computational cost is another point raised by this AI reviewer I found particularly important. Finally, this AI reviewer discusses issues with the code which I find very good, as a human reviewer I check the code only if I am very suspicious about the results and could have overlooked it. \\[3pt]
P7 $\cdot$ Claude 4.5 $\cdot$ paper-level slot 2 $\cdot$ primary & Yes there was a mention how systematic uncertainties were reported in Table I that was not spotted by other reviewers. \\[3pt]
P15 $\cdot$ Gemini 3.0 $\cdot$ paper-level slot 3 $\cdot$ secondary & [English translation from Korean] Clearly distinguishes sleep-structure abnormalities (local sleep disruption) from circadian-clock dysfunction. \\[3pt]
P23 $\cdot$ Gemini 3.0 $\cdot$ paper-level slot 1 $\cdot$ primary & It did significantly better for the depth of the comments, regarding the other 2 AI reviewers, but also made good point for some of the items, that were missed by the human reviewrs (frontal alpha in propofol). \\[3pt]
P24 $\cdot$ Claude 4.5 $\cdot$ paper-level slot 1 $\cdot$ primary & nice catching the normative atlas issue missed by the humans. \\[3pt]
P52 $\cdot$ Gemini 3.0 $\cdot$ paper-level slot 2 $\cdot$ primary & Yes. Raised issue of biosafety of gold dissolution. \\[3pt]
P52 $\cdot$ Gemini 3.0 $\cdot$ paper-level slot 2 $\cdot$ secondary & dissolution of gold for drug release can be toxic - needs testing \\[3pt]
P73 $\cdot$ Claude 4.5 $\cdot$ paper-level slot 2 $\cdot$ primary & AI Reviewer 2 was the only one to mention the potential confounding effect of population structure on the GWAS results. \\[3pt]
P73 $\cdot$ GPT-5.2 $\cdot$ paper-level slot 3 $\cdot$ primary & AI Reviewer 3 was the only one to raise a concern about the authors' extrapolation to estimate genome-wide mutation rates per microsatellite, although I consider this to be a reliatvely minor concern. \\[3pt]
P74 $\cdot$ GPT-5.2 $\cdot$ paper-level slot 2 $\cdot$ primary & AI Reviewer 2 noted that the non-human benchmark data were limited to a single SV caller and thus could have many errors. \\[3pt]
P77 $\cdot$ Claude 4.5 $\cdot$ paper-level slot 2 $\cdot$ primary & Yes, AI reviewer 2 pointed out am important point regarding the difference of Ag catalysis and Cu catalysis. \\[3pt]
P77 $\cdot$ GPT-5.2 $\cdot$ paper-level slot 3 $\cdot$ primary & AI reviewer 3 nicely brought up the issue of diastereoselectivity in the product formation which was not mentioned by the human reviewers. \\[3pt]
P81 $\cdot$ Gemini 3.0 $\cdot$ paper-level slot 1 $\cdot$ primary & - Only reviewer to question the use of the term ''LED,'' given that the device operates via gate oxide breakdown rather than conventional carrier injection. - Uniquely criticized the use of spatial intensity (mW/cm$^{2}$) as misleading, since the absolute EQE is extremely low compared to other Si emitters. \\[3pt]
P81 $\cdot$ Gemini 3.0 $\cdot$ paper-level slot 1 $\cdot$ secondary & AI Reviewer 1 was the only reviewer to question about the name ''LED''. It focused on how the light emission mechanism is different from traditional LEDs while other reviewers simply accepted the term. Second, while other reviewers also raised EQE concerns, it uniquely criticized the author's choice of the benchmark itself pointing out that the spatial intensity argument is misleading since the EQE is very small compared to other works. \\[3pt]
P83 $\cdot$ Gemini 3.0 $\cdot$ paper-level slot 1 $\cdot$ primary & - Only review to flag that 50 bending cycles at 5 mm radius is insufficient to justify the ''ultraflexible''. \\[3pt]
P83 $\cdot$ Gemini 3.0 $\cdot$ paper-level slot 1 $\cdot$ secondary & AI reviewer 1 focused on the mechanical durability compared to the standards in the field. \\[3pt]
P84 $\cdot$ GPT-5.2 $\cdot$ paper-level slot 2 $\cdot$ primary & - Only review to catch EM simulations used air/copper properties instead of tissue-equivalent properties. - Mentioned about insufficient strain sensor calibration based on fractional shortening metric. \\[3pt]
P84 $\cdot$ GPT-5.2 $\cdot$ paper-level slot 2 $\cdot$ secondary & AI reviewer 2 points out the insufficient calibration on the strain sensors, which makes cardiac contractility not convincing. \\[3pt]
\end{longtable}
}

\subsection*{S4: Internal consistency across sections --- n = 15}

{\footnotesize
\begin{longtable}{@{}p{3.2cm}p{11.5cm}@{}}
\toprule
\textbf{Citation} & \textbf{Expert comment} \\
\midrule
\endfirsthead
\toprule
\textbf{Citation} & \textbf{Expert comment} \\
\midrule
\endhead
\bottomrule
\endlastfoot
P1 $\cdot$ Gemini 3.0 $\cdot$ item 2 $\cdot$ primary & The point is that the goal may be to get closer to this accuracy and not to outperform it but at a reduced computational cost. The real critique here is that the authors don't clearly demonstrate how more and whether their approach is more computationally efficient. \\[3pt]
P37 $\cdot$ Claude 4.5 $\cdot$ item 1 $\cdot$ secondary & I wholeheartedly agree with this comment. This alone is enough to reject the paper, as the core claim doesn't match the empirical results. \\[3pt]
P50 $\cdot$ Claude 4.5 $\cdot$ item 3 $\cdot$ primary & Mistakes in writing - they switched the numbers. Numbers are correctly shown in Fig. 4f. \\[3pt]
P60 $\cdot$ Claude 4.5 $\cdot$ item 4 $\cdot$ secondary & This is a factual correction, as also noted by human reviewer 1 who caught a different labelling error in Item 4, this paper have slight proof reading issues. \\[3pt]
P60 $\cdot$ GPT-5.2 $\cdot$ item 3 $\cdot$ secondary & AI Reviewer 1's Item 5 noted the comparison is logically trivial (adding COVID deaths guarantees an increase); this reviewer shows that even setting that aside, the comparison is technically flawed in its construction. The inconsistent event definitions across periods is a particularly important point that no other reviewer raised. \\[3pt]
P61 $\cdot$ GPT-5.2 $\cdot$ item 5 $\cdot$ secondary & This is one of the most damaging critiques across all reviews because it challenges the paper's internal consistency rather than requiring external knowledge or judgment. A reader can verify every claim in this item solely from the manuscript. \\[3pt]
P65 $\cdot$ GPT-5.2 $\cdot$ item 2 $\cdot$ primary & This was one of my concerns regarding the original manuscript. They design for the 8-12 um range, but they their laser covers the 8.23 - 10.93 m range. \\[3pt]
P37 $\cdot$ Claude 4.5 $\cdot$ paper-level slot 2 $\cdot$ secondary & AI Reviewer 2 was the best in terms of analyzing what the authors claimed, and if the experiment design/results support the claims. \\[3pt]
P41 $\cdot$ Claude 4.5 $\cdot$ paper-level slot 3 $\cdot$ primary & AI Reviewer 3 caught the temporal mismatch issue. \\[3pt]
P41 $\cdot$ Gemini 3.0 $\cdot$ paper-level slot 1 $\cdot$ primary & AI Reviewer 1 caught the inconsistent projection methodology. \\[3pt]
P50 $\cdot$ Claude 4.5 $\cdot$ paper-level slot 3 $\cdot$ primary & Yes, AI Reviewer 3 found about mistakes in writing where authors reported numbers in wrong order and thus physiologically impossible number was reported. \\[3pt]
P50 $\cdot$ Claude 4.5 $\cdot$ paper-level slot 3 $\cdot$ secondary & AI Reviewer 3 correctly pointed out a discrepancy in the reported heart rate value - which was physiologically impossible for a healthy adult. \\[3pt]
P67 $\cdot$ GPT-5.2 $\cdot$ paper-level slot 2 $\cdot$ primary & Inconsistency in the number of epochs. \\[3pt]
P83 $\cdot$ GPT-5.2 $\cdot$ paper-level slot 3 $\cdot$ primary & - Uniquely caught that the abstract claims ''no passivation'' while Figure 2f's caption explicitly states the ultraflexible device used ''parylene encapsulation'' \\[3pt]
P83 $\cdot$ GPT-5.2 $\cdot$ paper-level slot 3 $\cdot$ secondary & AI reviewer 3 was the only reviewer to question about the encapsulation contradiction of ''no passivation'' and ''parylene encapsuation''. \\[3pt]
\end{longtable}
}

\subsection*{S5: Reproducibility / dependency failures --- n = 10}

{\footnotesize
\begin{longtable}{@{}p{3.2cm}p{11.5cm}@{}}
\toprule
\textbf{Citation} & \textbf{Expert comment} \\
\midrule
\endfirsthead
\toprule
\textbf{Citation} & \textbf{Expert comment} \\
\midrule
\endhead
\bottomrule
\endlastfoot
P54 $\cdot$ GPT-5.2 $\cdot$ item 4 $\cdot$ primary & significant point well supported that applies for reproducibility \\[3pt]
P66 $\cdot$ Gemini 3.0 $\cdot$ item 2 $\cdot$ primary & They do not provide the exact measurement results in their code, but one could extract the average data from the provided figures. \\[3pt]
P85 $\cdot$ GPT-5.2 $\cdot$ item 5 $\cdot$ primary & It is likely that upon further review the journal would require the deposition of these materials. \\[3pt]
P13 $\cdot$ Claude 4.5 $\cdot$ paper-level slot 2 $\cdot$ primary & Commented on the code availability; demographic mismatch \\[3pt]
P13 $\cdot$ GPT-5.2 $\cdot$ paper-level slot 3 $\cdot$ primary & Commented on the code availability; demographic mismatch \\[3pt]
P13 $\cdot$ Gemini 3.0 $\cdot$ paper-level slot 1 $\cdot$ primary & Commented on the code availability \\[3pt]
P42 $\cdot$ GPT-5.2 $\cdot$ paper-level slot 3 $\cdot$ primary & AI reviewer 3 mostly found reproducibility issues, which while valid, likely just require tidying and are highly unlikely to reflect scientific validity. \\[3pt]
P47 $\cdot$ GPT-5.2 $\cdot$ paper-level slot 2 $\cdot$ primary & AI \#2 identified the reproducibility failures in the experimental pipeline such as missing parameter specifications and modified external tools. \\[3pt]
P63 $\cdot$ GPT-5.2 $\cdot$ paper-level slot 3 $\cdot$ primary & Reproduccability item has not been raised by other AI and human reviewers. \\[3pt]
P68 $\cdot$ GPT-5.2 $\cdot$ paper-level slot 3 $\cdot$ primary & The code reproducibility issue was raised by AI reviewer 3 and not by others. \\[3pt]
\end{longtable}
}

\subsection*{S6: Big-picture / counter-narrative synthesis --- n = 7}

{\footnotesize
\begin{longtable}{@{}p{3.2cm}p{11.5cm}@{}}
\toprule
\textbf{Citation} & \textbf{Expert comment} \\
\midrule
\endfirsthead
\toprule
\textbf{Citation} & \textbf{Expert comment} \\
\midrule
\endhead
\bottomrule
\endlastfoot
P60 $\cdot$ Gemini 3.0 $\cdot$ item 4 $\cdot$ secondary & This is one of the most insightful critiques across all reviews because it runs counter to the paper's narrative rather than supporting the claim that compound impacts are underestimated, it suggests temperature-related mortality during the pandemic may actually be overestimated due to a less vulnerable surviving population. \\[3pt]
P85 $\cdot$ GPT-5.2 $\cdot$ item 2 $\cdot$ primary & Another way of phrasing this is that the authors over-interpret the connection between chromosome organization and metabolic state. \\[3pt]
P6 $\cdot$ Gemini 3.0 $\cdot$ paper-level slot 2 $\cdot$ primary & No, but this is primarily because it did not get bogged down in details. It was very good at seeing the big picture and pinpointing the biggest problems. \\[3pt]
P26 $\cdot$ Claude 4.5 $\cdot$ paper-level slot 1 $\cdot$ primary & this provided the most nuance in the comments, though all 3 AI did good job on noticing and calling up the circularity in the work, which was missed even by the first reviewer which was quite good. \\[3pt]
P26 $\cdot$ Gemini 3.0 $\cdot$ paper-level slot 3 $\cdot$ primary & good job on noticing and calling up the circularity in the work \\[3pt]
P37 $\cdot$ Gemini 3.0 $\cdot$ paper-level slot 1 $\cdot$ secondary & AI Reviewer 1 was the only reviewer who correctly noticed: 1. The training objective was not as novel as the authors claimed. 2. Insufficient comparison with state-of-the-art baselines. \\[3pt]
P49 $\cdot$ Gemini 3.0 $\cdot$ paper-level slot 3 $\cdot$ secondary & While other reviewers also expressed concerns about using connectome and molecular maps derived from healthy adults to explain the damage patterns of brain disorders, this AI reviewer 3 went a step further, demonstrating a deeper understanding of disorder-specific pathology. The logic is highly compelling: because neurodevelopmental disorders such as ASD and ADHD emerge at an early stage of life, unlike other brain conditions, their brain architecture fundamentally deviates even more from the normative connectomes and molecular maps of a healthy adult. Therefore, since these developmental disorders are the most likely to deviate from the normative baseline, it would significantly strengthen the study if the authors could obtain patient-specific data at least for these specific conditions to validate whether their observations still hold true. \\[3pt]
\end{longtable}
}

\subsection*{S\_generic: Residual: AI judged Correct without specific reason --- n = 40}

{\footnotesize
\begin{longtable}{@{}p{3.2cm}p{11.5cm}@{}}
\toprule
\textbf{Citation} & \textbf{Expert comment} \\
\midrule
\endfirsthead
\toprule
\textbf{Citation} & \textbf{Expert comment} \\
\midrule
\endhead
\bottomrule
\endlastfoot
P1 $\cdot$ Gemini 3.0 $\cdot$ item 3 $\cdot$ primary & I would also willing to see how their method performs with a larger cell. \\[3pt]
P6 $\cdot$ Claude 4.5 $\cdot$ item 3 $\cdot$ primary & This is well-argued and correct! \\[3pt]
P6 $\cdot$ Gemini 3.0 $\cdot$ item 2 $\cdot$ primary & Good point and well-argued! \\[3pt]
P7 $\cdot$ Gemini 3.0 $\cdot$ item 5 $\cdot$ primary & Yes this was mentioned previously as well, the paper should be rephrased a bit on this point, however the result is still valuable \\[3pt]
P9 $\cdot$ GPT-5.2 $\cdot$ item 3 $\cdot$ primary & What is being measured in Lb -\textgreater{} RR is not enough for an observation therefore the text is fine. \\[3pt]
P24 $\cdot$ GPT-5.2 $\cdot$ item 1 $\cdot$ primary & it is worth discussing it, I wouldn't say it is correct or incorrect. \\[3pt]
P25 $\cdot$ Gemini 3.0 $\cdot$ item 2 $\cdot$ primary & More references are needed. Roberts et al Nat Commun 2019 for example. \\[3pt]
P29 $\cdot$ Claude 4.5 $\cdot$ item 3 $\cdot$ primary & This is a good point. \\[3pt]
P29 $\cdot$ Claude 4.5 $\cdot$ item 5 $\cdot$ primary & Very good point. very well stated. \\[3pt]
P29 $\cdot$ Gemini 3.0 $\cdot$ item 4 $\cdot$ primary & This is quite good critique. \\[3pt]
P31 $\cdot$ Gemini 3.0 $\cdot$ item 2 $\cdot$ primary & Great point. \\[3pt]
P31 $\cdot$ Gemini 3.0 $\cdot$ item 4 $\cdot$ primary & Great point. \\[3pt]
P37 $\cdot$ Gemini 3.0 $\cdot$ item 3 $\cdot$ secondary & I wrote the paper ''Doctor AI'', and it is true what the AI reviewer is claiming here. \\[3pt]
P42 $\cdot$ Claude 4.5 $\cdot$ item 4 $\cdot$ primary & The AI is correct, while difficult to address methodologically, this is something that should be addressed in the discussion. \\[3pt]
P49 $\cdot$ GPT-5.2 $\cdot$ item 1 $\cdot$ secondary & By examining the materials used by authors more thoroughly than human reviewers, this reviewer was able to point out very important points. \\[3pt]
P53 $\cdot$ GPT-5.2 $\cdot$ item 1 $\cdot$ primary & very well supported observation \\[3pt]
P54 $\cdot$ Claude 4.5 $\cdot$ item 4 $\cdot$ primary & it catches detailed text that the authors should consider rewriting for clarity. \\[3pt]
P55 $\cdot$ Claude 4.5 $\cdot$ item 4 $\cdot$ primary & it is OK \\[3pt]
P61 $\cdot$ Claude 4.5 $\cdot$ item 1 $\cdot$ secondary & This is a major issue which none of the Human reviewers pointed but at the same time they all did emphasize more clarity on this figure especially. \\[3pt]
P61 $\cdot$ Gemini 3.0 $\cdot$ item 3 $\cdot$ secondary & This connects to Human Reviewer 1 Item 13 about Wang et al.: multiple prior studies did multi-factor analysis, and the paper should engage with them honestly rather than dismissing the entire prior literature. \\[3pt]
P62 $\cdot$ Claude 4.5 $\cdot$ item 1 $\cdot$ secondary & Addressing the same key point as Human reviewer \#2 and \#3, but the distinction between ex vivo and in vivo experiments were not detailed as in the comments from the Human reviewers. \\[3pt]
P64 $\cdot$ Gemini 3.0 $\cdot$ item 2 $\cdot$ primary & Good catch as AI reviewer 2 \\[3pt]
P65 $\cdot$ Claude 4.5 $\cdot$ item 3 $\cdot$ primary & I agree. A comparison of numerical vs. experimental Strehl ratios would be useful. \\[3pt]
P65 $\cdot$ Claude 4.5 $\cdot$ item 5 $\cdot$ primary & I am not sure how the authors could quantify the effect of LPA \\[3pt]
P68 $\cdot$ Claude 4.5 $\cdot$ item 3 $\cdot$ primary & Agree and also raised by the human reviewer 3 \\[3pt]
P68 $\cdot$ Claude 4.5 $\cdot$ item 4 $\cdot$ primary & This is an important point, especially when considering sexual dimorphism, geographic variation, etc. \\[3pt]
P68 $\cdot$ Gemini 3.0 $\cdot$ item 4 $\cdot$ primary & Some of the circularity is broken by the fact that analogs are defined based on multiple characteristics (not only inner ear). But the criticism is still relevant. \\[3pt]
P73 $\cdot$ Claude 4.5 $\cdot$ item 4 $\cdot$ primary & Again, I don't think that mutations in this gene must necessarily increase cancer risk to be candidates for slight increases in microsat mutation rates, but this would be worth discussing as it does weaken the case for this mutation somewhat. \\[3pt]
P74 $\cdot$ Claude 4.5 $\cdot$ item 4 $\cdot$ primary & I agree that this would have been useful to include. \\[3pt]
P77 $\cdot$ Claude 4.5 $\cdot$ item 5 $\cdot$ primary & This is a good point that was missed by the human reviewers. \\[3pt]
P77 $\cdot$ GPT-5.2 $\cdot$ item 1 $\cdot$ primary & The proposed 8-pi electrocyclization is assumed. \\[3pt]
P77 $\cdot$ GPT-5.2 $\cdot$ item 2 $\cdot$ secondary & The term ``stereo-controlled'' appears to be misleading. \\[3pt]
P85 $\cdot$ GPT-5.2 $\cdot$ item 1 $\cdot$ primary & While there are other samples C4 and C7 that are late time points in which data were generated with the same library kit, this is a valid concern. \\[3pt]
P85 $\cdot$ GPT-5.2 $\cdot$ item 4 $\cdot$ primary & I am not an expert on the analysis of these data, but this critique seems valid to me. \\[3pt]
P85 $\cdot$ Gemini 3.0 $\cdot$ item 1 $\cdot$ primary & I am not an expert in this analysis procedure, but this critique seems reasonable. \\[3pt]
P85 $\cdot$ Gemini 3.0 $\cdot$ item 2 $\cdot$ primary & I am not an expert in this analysis procedure, but this critique seems reasonable. \\[3pt]
P85 $\cdot$ Gemini 3.0 $\cdot$ item 3 $\cdot$ primary & I am not an expert in this analysis procedure, but this critique seems reasonable. \\[3pt]
P5 $\cdot$ Claude 4.5 $\cdot$ paper-level slot 3 $\cdot$ primary & Similar to my comments regarding the second AI reviewer, though the third AI reviewer provides stronger arguments. \\[3pt]
P35 $\cdot$ GPT-5.2 $\cdot$ paper-level slot 2 $\cdot$ secondary & Yes, All the Items. \\[3pt]
P67 $\cdot$ GPT-5.2 $\cdot$ paper-level slot 2 $\cdot$ secondary & I would say that the review of this AI is better than all the human reviews but not as good as the other AI reviews. Overall I find that the AI reviews are very good at judging the technical matters of the paper and human reviews have points that can only come from years of experience although not as thorough as the AI reviews. \\[3pt]
\end{longtable}
}
